% Дипломна работа, ОКС „Магистър“, специалност
% „Киберсигурност и превенция на киберпрестъпления“, ФКСТ, ТУ София.
% Preamble reused verbatim from the B.Eng thesis (doc/v2/bg/Main.tex),
% so that this document follows a format the faculty has already accepted.
% Do not edit here without reason; additions belong in Main.tex.
%
\documentclass[12pt]{article}

\usepackage{cmap}
\usepackage[T2A]{fontenc}
\usepackage[utf8]{inputenc}
\usepackage[bulgarian]{babel}

% --- Типография ---
\usepackage{tempora}
\usepackage[top=2.5cm, bottom=2.5cm, left=3.0cm, right=2.0cm, a4paper]{geometry}
\usepackage{setspace}
\setstretch{1.5}
\usepackage[none]{hyphenat}
\usepackage[expansion=false]{microtype}
\DisableLigatures{encoding=*,family=tt*}
\emergencystretch=1.5em
\sloppy
\usepackage{makecell}
\usepackage{indentfirst}
\setlength{\parindent}{1.25cm}
\setlength{\parskip}{0pt}

% --- Номерация на страниците (долу вдясно) ---
\usepackage{fancyhdr}
\pagestyle{fancy}
\fancyhf{}
\renewcommand{\headrulewidth}{0pt}
\fancyfoot[R]{\thepage}

% --- Бележки под линия ---
\usepackage[hang,flushmargin]{footmisc}
\usepackage{ragged2e}
\renewcommand{\footnotelayout}{\fontsize{10}{12}\selectfont\setstretch{1.0}\justifying}

% --- Форматиране на заглавия ---
\usepackage{titlesec}
\newcommand{\sectionbreak}{\clearpage}

\renewcommand{\thesection}{\Roman{section}}
\renewcommand{\thesubsection}{\arabic{section}.\arabic{subsection}}
\renewcommand{\thesubsubsection}{\arabic{section}.\arabic{subsection}.\arabic{subsubsection}}

% Sections: Roman, Bold, 14pt, Centered, Uppercase
\titleformat{\section}[block]
  {\fontsize{14}{16.8}\bfseries\centering}
  {\thesection. }
  {0pt}
  {\MakeUppercase}

% Paragraphs: Arabic, Bold, 14pt, 1.25cm indent
\titleformat{\subsection}[block]
  {\fontsize{14}{16.8}\bfseries}
  {\thesubsection. }
  {0pt}
  {}
\titlespacing*{\subsection}{1.25cm}{*3}{*2}

% Sub-paragraphs: Bold, 13pt, 2.0cm indent
\titleformat{\subsubsection}[block]
  {\fontsize{13}{15.6}\bfseries}
  {\thesubsubsection. }
  {0pt}
  {}
\titlespacing*{\subsubsection}{2.0cm}{*2}{*1}

% --- Съдържание ---
\usepackage{etoolbox}
\makeatletter
\patchcmd{\l@section}{#1}{\MakeUppercase{#1}}{}{}
\makeatother
\usepackage[titles]{tocloft}
\renewcommand{\cftsecleader}{\cftdotfill{\cftdotsep}}
\renewcommand{\cftsubsecleader}{\cftdotfill{\cftdotsep}}
\renewcommand{\cftsubsubsecleader}{\cftdotfill{\cftdotsep}}
\renewcommand{\cftsecfont}{\fontsize{12}{14}\bfseries}
\renewcommand{\cftsecpagefont}{\fontsize{12}{14}\bfseries}
\renewcommand{\cftsecaftersnum}{.\enspace}
\settowidth{\cftsecnumwidth}{VIII.\enspace}
\renewcommand{\cftsubsecaftersnum}{.}
\renewcommand{\cftsubsubsecaftersnum}{.}

% --- Надписи (таблици и фигури) ---
\usepackage{caption}
\usepackage{subcaption}
\DeclareCaptionLabelFormat{tablelabel}{Таблица #2}
\captionsetup[table]{
    labelformat=tablelabel,
    labelsep=period,
    font={bf,small},
    justification=centering,
    singlelinecheck=true
}
\DeclareCaptionLabelFormat{figurelabel}{Фигура #2}
\captionsetup[figure]{
    labelformat=figurelabel,
    labelsep=period,
    font={bf,small},
    justification=centering,
    singlelinecheck=true
}

% --- Математика ---
\usepackage{amsmath,amssymb,amsfonts}
\usepackage{amsthm}

% --- Графики ---
\usepackage{graphicx}
\usepackage{float}
\usepackage{tikz}
\usetikzlibrary{decorations.pathreplacing,arrows.meta,positioning,shapes.geometric,fit,calc}
\usepackage{adjustbox}

% --- Таблици ---
\usepackage{array}
\usepackage{booktabs}
\usepackage{longtable}
\usepackage{multirow}
\usepackage{tabularx}
% arydshln is NOT loaded. Under the LaTeX kernel of 2026-06-01 it breaks every
% tabularx in the document with "Extra }, or forgotten $" (measured: the
% unmodified B.Eng thesis fails the same way on its first table). Dashed rules
% fall back to a plain booktabs rule.
\providecommand{\hdashline}{\midrule}
% (arydshln not loaded, see above) \setlength\dashlinedash{1.5pt}
% (arydshln not loaded, see above) \setlength\dashlinegap{1.5pt}
\setlength\arrayrulewidth{0.4pt}

% --- Библиография ---
\usepackage{hyperref}
\usepackage{caption}

% --- Помощни пакети ---
\usepackage{xcolor}
\usepackage{xurl}
\usepackage{soul}
\usepackage{enumerate}
\usepackage{enumitem}
\usepackage{textcomp}
\usepackage{gensymb}
\usepackage{fancyvrb}
\usepackage{eurosym}
\usepackage{ragged2e}
\usepackage{siunitx}
\usepackage{listings}
\usepackage{seqsplit}
\usepackage{collcell}

\captionsetup{font=small,labelfont=bf}
\urlstyle{same}
\newcommand{\ORMName}{Compile-time Object-Relational Mapping}
\pdfstringdefDisableCommands{\def\code#1{#1}\def\texttt#1{#1}\def\ORMName{Compile-time Object-Relational Mapping}}
\ExplSyntaxOn
\tl_new:N \l__orm_code_text_tl
\cs_new_protected:Npn \orm_code:n #1
  {
    \group_begin:
      \tl_set:Nx \l__orm_code_text_tl { \tl_to_str:n {#1} }
      \regex_replace_all:nnN { \\([_&]) } { \1 } \l__orm_code_text_tl
      \regex_replace_all:nnN { \\\{ } { \c{textbraceleft} } \l__orm_code_text_tl
      \regex_replace_all:nnN { \\\} } { \c{textbraceright} } \l__orm_code_text_tl
      \regex_replace_all:nnN { < } { \c{textless} } \l__orm_code_text_tl
      \regex_replace_all:nnN { > } { \c{textgreater} } \l__orm_code_text_tl
      \mbox { \ttfamily \tl_use:N \l__orm_code_text_tl }
    \group_end:
  }
\DeclareRobustCommand{\code}[1]{\orm_code:n {#1}}
\cs_new_protected:Npn \orm_inline_text:n #1
  {
    \regex_match:nnTF { \A C\+\+(?:\d+)? \Z } {#1}
      {#1}
      {
        \regex_match:nnTF { [<_:\\.] } {#1}
          {\code{#1}}
          {
            \regex_match:nnTF { \A [-A-Z0-9/+_]+(?:\s+[-A-Z0-9/+_]+)+ \Z } {#1}
              {\code{#1}}
              {
                \tl_if_in:nnTF {#1} { }
                  {#1}
                  {\code{#1}}
              }
          }
      }
  }
\DeclareRobustCommand{\texttt}[1]{\orm_inline_text:n {#1}}
\ExplSyntaxOff
\setlength{\emergencystretch}{3em}
\newcommand{\ORMTableCellStyle}{\RaggedRight\hspace{0pt}\relax}
\ExplSyntaxOn
\cs_new_protected:Npn \orm_insert_slash_breaks:n #1
  {
    \tl_set:Nn \l_tmpa_tl { #1 }
    \tl_replace_all:Nnn \l_tmpa_tl { / } { /\allowbreak }
    \tl_use:N \l_tmpa_tl
  }
\NewDocumentCommand{\ORMSlashTableCell}{m}
  {{\ORMTableCellStyle \orm_insert_slash_breaks:n {#1}}}
\ExplSyntaxOff
\newcolumntype{L}[1]{>{\ORMTableCellStyle\arraybackslash}p{#1}}
\newcolumntype{Y}{>{\ORMTableCellStyle\arraybackslash}X}
\newcolumntype{Q}[1]{>{\collectcell\ORMSlashTableCell}p{#1}<{\endcollectcell}}
\newcolumntype{Z}{>{\collectcell\ORMSlashTableCell}X<{\endcollectcell}}
\setlength{\extrarowheight}{0.45em}
\tikzset{
    ormblock/.style={draw, rounded corners, fill=blue!5, align=center, font=\small,
        minimum width=3.2cm, minimum height=1.05cm, inner sep=6pt},
    ormaccent/.style={draw, rounded corners, fill=green!8, align=center, font=\small,
        minimum width=3.0cm, minimum height=0.95cm, inner sep=6pt},
    ormdanger/.style={draw, rounded corners, fill=red!6, align=center, font=\small,
        minimum width=3.0cm, minimum height=0.95cm, inner sep=6pt},
    ormnote/.style={draw, dashed, rounded corners, fill=gray!8, align=left, font=\small,
        inner sep=6pt},
    ormline/.style={-Latex, thick}
}

% --- Единично разстояние в нетекстови среди ---
\AtBeginEnvironment{tikzpicture}{\setstretch{1.1}}
\AtBeginEnvironment{lstlisting}{\setstretch{1.1}}
\AtBeginEnvironment{verbatim}{\setstretch{1.1}}
\AtBeginEnvironment{tabular}{\setstretch{1.1}}
\AtBeginEnvironment{longtable}{\setstretch{1.1}}

\BeforeBeginEnvironment{tikzpicture}{\begin{adjustbox}{max width=\textwidth}}
\AfterEndEnvironment{tikzpicture}{\end{adjustbox}}

% --- Оцветяване на код ---
\definecolor{cppblue}{RGB}{0,0,180}
\definecolor{cppgreen}{RGB}{0,128,0}
\definecolor{cppred}{RGB}{180,0,0}
\definecolor{cppgray}{RGB}{128,128,128}

\usepackage[T2A]{fontenc}
\lstset{
    language=C++,
    basicstyle=\footnotesize\ttfamily,
    breaklines=true,
    breakatwhitespace=false,
    breakautoindent=false,
    breakindent=0pt,
    columns=fullflexible,
    keepspaces=true,
    postbreak=\mbox{\textcolor{cppgray}{\ensuremath{\hookrightarrow}}\space},
    xleftmargin=1.5em,
    xrightmargin=0.5em,
    keywordstyle=\color{cppblue}\bfseries,
    stringstyle=\color{cppred},
    commentstyle=\color{cppgreen},
    identifierstyle=\color{black},
    showstringspaces=false,
    numbers=left,
    numberstyle=\tiny\color{cppgray},
    numbersep=8pt,
    aboveskip=1.5em,
    belowskip=1.5em,
    frame=single,
    framerule=0.5pt,
    rulecolor=\color{cppgray},
    tabsize=4,
    captionpos=b,
    morekeywords={co_await, co_return, co_yield, Task, expected, unexpected,
                  requires, concept, constexpr, noexcept, nullptr, override, final,
                  ph, prepare, execute, select, insert, update, deleteq,
                  Placeholder, IndexedPlaceholder, field, where, order_by, group_by, limit},
    emph={orm, db, result, prepared_query, connector_trait, select_query,
          insert_query, update_query, delete_query, MockDB, SQLiteDB},
    emphstyle=\color{cppblue},
    inputencoding=utf8,
    literate={а}{{\cyra}}1 {б}{{\cyrb}}1 {в}{{\cyrv}}1 {г}{{\cyrg}}1 {д}{{\cyrd}}1
             {е}{{\cyre}}1 {ж}{{\cyrzh}}1 {з}{{\cyrz}}1 {и}{{\cyri}}1 {й}{{\cyrishrt}}1
             {к}{{\cyrk}}1 {л}{{\cyrl}}1 {м}{{\cyrm}}1 {н}{{\cyrn}}1 {о}{{\cyro}}1
             {п}{{\cyrp}}1 {р}{{\cyrr}}1 {с}{{\cyrs}}1 {т}{{\cyrt}}1 {у}{{\cyru}}1
             {ф}{{\cyrf}}1 {х}{{\cyrh}}1 {ц}{{\cyrc}}1 {ч}{{\cyrch}}1 {ш}{{\cyrsh}}1
             {щ}{{\cyrshch}}1 {ъ}{{\cyrhrdsn}}1 {ь}{{\cyrsftsn}}1 {ю}{{\cyryu}}1 {я}{{\cyrya}}1
             {А}{{\CYRA}}1 {Б}{{\CYRB}}1 {В}{{\CYRV}}1 {Г}{{\CYRG}}1 {Д}{{\CYRD}}1
             {Е}{{\CYRE}}1 {Ж}{{\CYRZH}}1 {З}{{\CYRZ}}1 {И}{{\CYRI}}1 {Й}{{\CYRISHRT}}1
             {К}{{\CYRK}}1 {Л}{{\CYRL}}1 {М}{{\CYRM}}1 {Н}{{\CYRN}}1 {О}{{\CYRO}}1
             {П}{{\CYRP}}1 {Р}{{\CYRR}}1 {С}{{\CYRS}}1 {Т}{{\CYRT}}1 {У}{{\CYRU}}1
             {Ф}{{\CYRF}}1 {Х}{{\CYRH}}1 {Ц}{{\CYRC}}1 {Ч}{{\CYRCH}}1 {Ш}{{\CYRSH}}1
             {Щ}{{\CYRSHCH}}1 {Ъ}{{\CYRHRDSN}}1 {Ь}{{\CYRSFTSN}}1 {Ю}{{\CYRYU}}1 {Я}{{\CYRYA}}1,
}


\providecommand{\TODO}[1]{\textcolor{red}{[TODO: #1]}}
\hypersetup{
    pdftitle={Предотвратяване на инжекционни атаки чрез конструиране на заявки по време на компилация},
    pdfauthor={Алекс Цветанов},
}

\begin{document}

% --- Заглавна страница ---
\begin{titlepage}
\begin{center}
    {\fontsize{18}{21.6}\selectfont\bfseries ТЕХНИЧЕСКИ УНИВЕРСИТЕТ \textendash{} СОФИЯ\par}
    \vspace{0.5cm}
    {\fontsize{16}{19.2}\selectfont\bfseries ФАКУЛТЕТ „КОМПЮТЪРНИ СИСТЕМИ И ТЕХНОЛОГИИ“\par}

    \vspace{0.5cm}
    \begin{figure}[H]
        \centering
        \includegraphics[width=0.4\linewidth]{images/ту.png}
    \end{figure}
    \vspace{0.5cm}

    {\fontsize{18}{21.6}\selectfont\bfseries ПРЕДОТВРАТЯВАНЕ НА ИНЖЕКЦИОННИ АТАКИ ЧРЕЗ КОНСТРУИРАНЕ НА ЗАЯВКИ ПО ВРЕМЕ НА КОМПИЛАЦИЯ: АНАЛИЗ И ОЦЕНКА ПРИ ХЕТЕРОГЕННИ ХРАНИЛИЩА ЗА ДАННИ\par}

    \vspace{1cm}

    {\fontsize{14}{16.8}\selectfont\bfseries ДИПЛОМНА РАБОТА\par}
    \vspace{0.3cm}
    {\fontsize{14}{16.8}\selectfont\bfseries за придобиване на ОКС „Магистър“\par}

    \vfill

    \begin{flushleft}
    {\fontsize{14}{16.8}\selectfont\bfseries Изготвил:\par}
    Алекс Цветанов, фак. № 371326005\par
    Специалност: Киберсигурност и превенция на киберпрестъпления\par

    \vspace{1cm}

    {\fontsize{14}{16.8}\selectfont\bfseries Научен ръководител:\par}
    \TODO{звание, име на ръководителя}\par
    \end{flushleft}

    \vfill

    {\fontsize{12}{14.4}\selectfont София, 2027\par}
\end{center}
\end{titlepage}

\newpage
\pagestyle{empty}
\section*{Резюме}

Инжекционните атаки остават сред най-разпространените уязвимости в софтуера и са
средство за редица документирани киберпрестъпления. Настоящата дипломна работа
изследва дали една абстракция, конструираща заявките към хранилища за данни по
време на компилация, може да елиминира цяло семейство от класове инжекции при
хетерогенни хранилища, и къде точно спира тази гаранция. Като платформа се
използва библиотека за обектно-релационно картографиране на C++23, разработена в
предходна дипломна работа, която работи със седем хранилища: PostgreSQL, MySQL,
SQLite, Cassandra, MongoDB, Neo4j и Redis.

Въз основа на определението на инжекцията като промяна на синтактичната
структура на заявката е формулиран аргумент: когато имената на таблици и колони и
формата на заявката са аргументи на шаблони, а стойностите достигат заявката само
чрез свързани типизирани заместители, недоверен вход не може да промени
структурата. Аргументът е проверен с три независими метода: 26 отрицателни
компилационни теста и 5 контролни случая, с проверка на диагностиката при GCC 14.2
и Clang 18.1; диференциален тест, при който всеки приложим полезен товар от корпус
от 66 публикувани полезни товара от 11 източника се подава към съзнателно уязвима
заявка и към библиотеката; и анализ на изходния код на всеки конектор.

Анализът установи, че гаранцията се постига с четири различни механизма, и откри
пропуск в нея: стойност, записана направо в заявката вместо чрез заместител, се
приемаше, но не се свързваше, и заявката тихо връщаше грешни резултати. Пропускът
е затворен с проверка по време на компилация във всички седем входни точки, при
запазено успешно изпълнение на всички съществуващи тестове. При SQLite
наивната заявка беше компрометирана от 4 от 22 приложими полезни товара,
включително с извличане на съхранена идентификационна информация, а библиотеката
от нито един. Тестът беше изпълнен и срещу работещи сървъри на PostgreSQL, MySQL,
MongoDB, Redis, Cassandra и Neo4j в система за непрекъсната интеграция. Общо за
седемте хранилища през библиотеката не е имала ефект нито една от 98-те
двойки на полезен товар и хранилище, съставени от всички 66 полезни товара.
Наивният път е компрометиран от по 2 полезни товара при PostgreSQL, MySQL,
MongoDB, Redis и Neo4j. При Cassandra той не е компрометиран от нито един и
сравнението е неубедително.

Библиотеката беше разширена в две посоки, в които гаранцията не е очевидна. При
условия \code{IN} и \code{NOT IN} вход по време на изпълнение определя дължината
на списък, но текстът на заявката получава само броя на маркерите. Миграциите
работят с имена, прочетени обратно от самото хранилище, затова там гаранцията е
проверка по време на изпълнение, а не компилация. Тя е измерена и при инжекция от
втори ред: наивният DDL е компрометиран при 14 от 368 изпълнения,
всички от съставени входове, а библиотеката при нито едно.

Приносите на работата са аргумент за предотвратяване на инжекциите чрез
конструиране по време на компилация, разширен извън SQL; възпроизводим набор за
оценка с корпус с посочени източници; открит и затворен пропуск в гаранцията; и
определена граница на гаранцията при списъци и операции по схемата.

\bigskip
\noindent\textbf{Ключови думи:} инжекционни атаки, предотвратяване на
киберпрестъпления, сигурност по проектиране, C++, програмиране по време на
компилация, хранилища за данни.

\section*{Abstract}

Injection attacks remain among the most common software vulnerabilities and are
the means of a number of documented cybercrimes. This thesis examines whether an
abstraction that constructs queries to data stores at compile time can eliminate
a whole family of injection classes across heterogeneous stores, and where
exactly that guarantee stops. The platform is a C++23 object-relational mapping
library, developed in an earlier thesis, that works with seven stores:
PostgreSQL, MySQL, SQLite, Cassandra, MongoDB, Neo4j and Redis.

Building on the definition of injection as a change to the syntactic structure
of a query, the thesis argues that when table names, column names and the shape
of the query are template arguments, and values reach the query only through
bound, typed placeholders, untrusted input cannot change that structure. The
argument is tested by three independent methods: 26 negative compilation tests
and 5 control cases, with diagnostic matching under GCC 14.2 and Clang 18.1; a
differential test that sends every applicable payload from a corpus of 66
published payloads from 11 sources both to a deliberately vulnerable query and to
the library; and an analysis of the source of every connector.

The analysis found that the guarantee is achieved by four different mechanisms,
and found a gap in it: a value written directly into a query instead of through a
placeholder was accepted but never bound, and the query silently returned wrong
results. The gap is closed by a compile-time check at all seven entry points,
with all existing tests still passing. On SQLite the vulnerable query
was compromised by 4 of 22 applicable payloads, including the theft of a stored
credential, and the library by none. The test was then run against live servers
of PostgreSQL, MySQL, MongoDB, Redis, Cassandra and Neo4j in a continuous
integration system. Across the seven stores, none of the 98 pairs of payload
and store, covering all 66 payloads, had an effect through the library. The naive
path was compromised by 2 payloads each on PostgreSQL, MySQL, MongoDB, Redis and
Neo4j. On Cassandra it was never compromised, so the comparison is inconclusive.

The library was extended in two directions where the guarantee is not obvious. For
\code{IN} and \code{NOT IN}, run-time input sets the length of a list, but the query
text receives only the number of markers. Migrations work with names read back from
the store itself, so there the guarantee is a run-time check and not a compile-time
one. It was measured, including second order injection: the naive DDL was compromised
in 14 of 368 runs, all by constructed inputs, and the library in none.

The contributions are an argument for preventing injection by compile-time
construction, extended beyond SQL; a reproducible evaluation suite with a sourced
corpus; a gap in the guarantee, found and closed; and the boundary of the guarantee
for lists and schema operations, located and measured.

\bigskip
\noindent\textbf{Keywords:} injection attacks, cybercrime prevention, secure by
design, C++, compile-time programming, data stores.


\newpage
\setcounter{tocdepth}{2}
\tableofcontents
\newpage

\pagestyle{fancy}
\pagenumbering{arabic}

\section*{УВОД}
\label{sec:introduction}
\addcontentsline{toc}{section}{УВОД}

\subsection*{Актуалност}

Инжекцията е една от най-старите и най-устойчивите уязвимости в софтуера. Тя
продължава да заема видимо място в класациите на най-сериозните рискове за
уеб приложенията~\cite{owasp_top10} и е средство за документирани
киберпрестъпления, обсъдени в глава~\ref{sec:cybercrime}. Устойчивостта ѝ не се
дължи на липса на знание. Защитата срещу класическата SQL инжекция е известна
отдавна: стойностите трябва да се подават като параметри, а не да се вграждат в
текста на заявката~\cite{owasp_sqli_prevention}. Уязвимостта оцелява, защото
защитата зависи от дисциплината на всеки разработчик при всяка заявка.

Това е причината подходът към инжекциите да се измества от откриване и поправяне
на отделни случаи към тяхното изключване като клас. Съвместно предупреждение на
американските агенции CISA и FBI призовава производителите на софтуер да
елиминират SQL инжекциите именно по този начин~\cite{cisa_sqli_alert}. Възниква
въпросът какви средства позволяват такова изключване и докъде то стига.

\subsection*{Изследователски въпрос}

Настоящата работа изследва един такъв подход: конструиране на заявките по време
на компилация. При него структурата на заявката, имената на таблиците и на
колоните са фиксирани от типовата система, а стойностите могат да влязат само
през типизирани заместители. Работата отговаря на въпроса:

\begin{quote}
\textit{Може ли една абстракция, конструираща заявките по време на компилация, да
елиминира цяло семейство от класове инжекции при хетерогенни хранилища за
данни, и къде точно спира тази гаранция?}
\end{quote}

Въпросът е поставен за семейство от класове, а не само за SQL, защото
съвременните приложения използват различни видове хранилища, всяко със собствен
език и собствен клас инжекции. Втората част на въпроса е равностойна на първата.
Гаранция, чиито граници не са известни, не може да бъде използвана безопасно.

\subsection*{Цел и задачи}

Целта на работата е да провери, чрез измерване, а не само чрез разсъждение,
доколко конструирането на заявки по време на компилация предотвратява
инжекционните атаки при седем хетерогенни хранилища за данни, и да установи
границите на тази защита. За постигане на целта са поставени следните задачи:

\begin{enumerate}
    \item да се анализира инжекцията като средство за киберпрестъпления и
          правната рамка за противодействието ѝ;
    \item да се прегледат съществуващите защити и да се установи незапълнената
          нужда;
    \item да се формулира аргумент, свързващ формалното определение на
          инжекцията~\cite{su2006essence} с типовата система на C++;
    \item да се определи модел на заплахите и критерии за успех, фиксирани преди
          измерванията;
    \item да се анализира реализацията на гаранцията при всяко от седемте
          хранилища;
    \item да се изгради възпроизводим набор за оценка: отрицателни компилационни
          тестове, корпус от публикувани полезни товари и диференциален тест;
    \item да се оценят резултатите спрямо критериите и да се назоват границите на
          гаранцията.
\end{enumerate}

\subsection*{Обхват}

Обхватът е пътят на заявката от изграждането ѝ в приложението до предаването ѝ
на драйвера на хранилището. Извън обхвата са операциите по схемата, инжекцията от
втори ред, логическите грешки в приложенията, уязвимостите в самите драйвери и
отказът от услуга чрез скъпи заявки. Тези граници са обосновани в
точка~\ref{sec:tm-scope}.

Работата използва като платформа библиотека за обектно-релационно
картографиране, разработена в предходна дипломна работа за ОКС
„Бакалавър“~\cite{tsvetanov2026orm}. Предходната работа описва проектирането и
реализацията на библиотеката и не разглежда въпроса за сигурността. Настоящата
работа не претендира за нейните резултати. Тя изследва самостоятелен въпрос и
формулира приноси, които се отнасят само към него и към промените в
библиотеката, произтекли от изследването.

\subsection*{Структура на работата}

Глава~\ref{sec:cybercrime} разглежда инжекцията като средство за
киберпрестъпления и правната рамка за противодействието ѝ.
Глава~\ref{sec:existing} прави преглед на съществуващите защити и установява
незапълнената нужда. Глава~\ref{sec:theory} въвежда формалното определение на
инжекцията и аргумента за времето на компилация като граница на доверие.
Глава~\ref{sec:threat-model} определя модела на заплахите, критериите за успех и
метода на оценката. Глава~\ref{sec:enforcement} анализира как гаранцията се
постига при седемте хранилища. Глава~\ref{sec:case-study} проследява открит в хода
на работата пропуск в гаранцията, от откриването до поправката.
Глава~\ref{sec:evaluation} представя и тълкува резултатите от оценката.
Заключението обобщава приносите, ограниченията и насоките за бъдеща работа.

\newpage
\section{Инжекциите като средство за киберпрестъпления}
\label{sec:cybercrime}

Уводът поставя въпроса дали един механизъм, действащ по време на компилация, може
да изключи цял клас инжекции. Преди да се търси отговор, трябва да се покаже защо
въпросът има значение извън инженерната практика. Настоящата глава разглежда
инжекцията като средство за престъпления и като източник на отговорност за онези,
които обработват чужди данни. Първо се очертава мястото на инжекциите в
класациите на уязвимостите. След това се разглеждат три документирани случая и
правната рамка: международна, европейска и българска. Накрая се разглежда преходът
от наказване към превенция, който мотивира работата.

\subsection{Инжекцията в съвременния пейзаж на заплахите}
\label{sec:cybercrime-landscape}

\textbf{Инжекция} в широк смисъл е ситуация, при която недоверен вход бъде
изтълкуван от интерпретатор като част от команда, а не като данни. Формалното
определение, което работата използва, се въвежда в глава~\ref{sec:theory}.

Действащото издание на класацията на \textit{OWASP} за най-сериозните рискове при
уеб приложенията е \textit{OWASP Top 10:2025}~\cite{owasp_top10}. В него
инжекцията е на пето място (категория A05:2025), след като е слязла с две места
от трето място в предишното издание. Категорията обединява 37 класа слабости по
номенклатурата \textit{CWE} и според документа е с най-голям брой регистрирани
уязвимости (\textit{CVE}) от всички категории~\cite{owasp_a05_injection}.

Номенклатурата \textit{CWE} (\textit{Common Weakness Enumeration}) на \textit{MITRE}
описва видовете слабости, а не отделните уязвимости. Записът CWE-89 е озаглавен
„Improper Neutralization of Special Elements used in an SQL Command ('SQL
Injection')“~\cite{cwe89}. Записът CWE-943 е озаглавен „Improper Neutralization of
Special Elements in Data Query Logic“ и обхваща заявка към хранилище на данни,
например база от данни, без да се ограничава до SQL~\cite{cwe943}. Тази разлика в
обхвата е съществена за работата. В класацията \textit{2025 CWE Top 25 Most
Dangerous Software Weaknesses} CWE-89 е на второ място, с едно място нагоре спрямо
2024 година~\cite{cwe_top25_2025}. Това е най-новият списък, наличен на страницата
на проекта към 4 ноември 2026 година.

Тези източници не доказват, че инжекциите се използват за престъпления. Те
показват, че класът е широко разпространен, устойчив във времето и добре описан,
т.е. че защитата срещу него е въпрос на познати средства, а не на неизвестна заплаха.

\subsection{Документирани случаи}
\label{sec:cybercrime-cases}

За всеки случай има първичен или официален източник, който го свързва със SQL
инжекция. Цифри са посочени само там, където са прочетени в такъв източник.
Случаите илюстрират три правни положения: масово изнудване с откраднати данни,
наказателно преследване на извършител и административна санкция на администратор
на данни.

\subsubsection{MOVEit Transfer и CL0P (2023)}

Уязвимостта \textit{CVE-2023-34362} засяга приложението \textit{MOVEit Transfer} на
\textit{Progress Software}. Записът в Националната база данни за уязвимости на
\textit{NIST} я описва като SQL инжекция, която може да позволи на неудостоверен
нападател достъп до базата от данни, класифицира я като CWE-89 и ѝ дава оценка 9.8
по \textit{CVSS v3.1}~\cite{nvd_cve_2023_34362}. Съвместното предупреждение
\textit{AA23-158A} на \textit{FBI} и \textit{CISA} от 7 юни 2023 година приписва кампанията на
групата \textit{CL0P}. Според него нападателите са използвали уязвимостта, за да
инсталират уеб обвивката \textit{LEMURLOOT} и да изтеглят данни от базите на
приложението, а от 2021 година групата разчита главно на кражба на данни, а не на
криптиране. Използването е започнало на 27 май 2023 година според публични
източници, а на 2 юни 2023 година уязвимостта е добавена в каталога на активно
използваните~\cite{cisa_aa23_158a}.

\subsubsection{Heartland Payment Systems и Албърт Гонсалес (2009, 2010)}

На 17 август 2009 година прокуратурата на щата Ню Джърси съобщава за обвинителен
акт срещу трима души за пет корпоративни пробива, сред които \textit{Heartland
Payment Systems}, \textit{7-Eleven} и \textit{Hannaford}. Обвинителният акт описва
кражба на над 130 милиона номера на кредитни и дебитни карти от тези три компании и
определя атаките като „SQL-injection attack“. Съобщението напомня, че
обвинителният акт е само обвинение~\cite{doj_gonzalez_indictment_2009}. На 26 март
2010 година Министерството на правосъдието на САЩ съобщава, че Албърт Гонсалес е
осъден на 20 години и един ден лишаване от свобода по две обвинения в съзаговор
във връзка с достъпа до платежните мрежи на тези три компании. Присъдата е
едновременна с постановената предния ден в Бостън по друго
дело~\cite{doj_gonzalez_sentence_2010}. Съобщението за присъдата не използва
термина „SQL инжекция“. Приписването на метода се основава на обвинителния акт, а не
на съдебно установяване.

\subsubsection{TalkTalk и санкцията на ICO (2016)}

На 5 октомври 2016 година британската Служба на комисаря по информацията
(\textit{Information Commissioner's Office}, \textit{ICO}) налага на
\textit{TalkTalk} глоба от 400\,000 британски лири, определена от службата като
рекордна. Според нея атаката е осъществена от 15 до 21 октомври 2015 година чрез
три уязвими уеб страници в наследена инфраструктура. Нападателят е получил достъп до
личните данни на 156\,959 клиенти, като при 15\,656 от тях е имал достъп и до данни
за банкови сметки~\cite{ico_talktalk_2016}. Службата установява, че нападателят е
използвал SQL инжекция, че техниката е „well understood“, защити съществуват и
компанията е трябвало да знае за риска. Имало е и две незабелязани предупреждения:
успешна SQL инжекция на 17 юли 2015 година срещу същата уязвимост и втора атака
между 2 и 3 септември 2015 година. Нарушението е квалифицирано като неспазване на
седмия принцип на британския \textit{Data Protection Act 1998}, а отделно
наказателно разследване води Столичната полиция~\cite{ico_talktalk_2016}.
Санкцията е наложена по действалия тогава закон, а не по Общия регламент относно
защитата на данните. Тя е за липса на мерки срещу добре позната атака, а не за
самата атака.

Трите случая заемат три различни роли: обект на масово изнудване, извършител,
подведен под наказателна отговорност, и администратор, санкциониран за липса на
защита. Общото е, че слабостта е от клас, за който защитата е известна отдавна.
\textit{ICO} го казва за конкретния случай, а предупреждението на \textit{CISA} и
\textit{FBI} (раздел~\ref{sec:cybercrime-prevention}), което говори за широко огласена кампания срещу приложение за управляван трансфер на файлове, засегнала
хиляди организации, го казва за индустрията. То не назовава продукта~\cite{cisa_sqli_alert};
връзката с \textit{MOVEit} е извод на автора. Проблемът не е липсата на знание, а
това, че знанието се прилага ръчно, на всяко място поотделно.

\subsection{Правна рамка}
\label{sec:cybercrime-law}

Правното третиране на инжекцията се определя от резултата, а не от техниката.
Нападател, който чрез инжекция получи достъп до данни, ги промени или спре работата
на система, може да осъществи състава на престъпление по международни, европейски и
национални актове. Администраторът на данните носи самостоятелно задължение да
осигури защита. Прилагането на състав към конкретно деяние е въпрос на съдебна
преценка, затова връзката между инжекцията и даден състав по-долу е формулирана като
възможна правна квалификация. Съдебна практика не се разглежда.

\subsubsection{Конвенцията от Будапеща}

Конвенцията за престъпленията в кибернетичното пространство на Съвета на Европа
(\textit{ETS № 185}, Будапеща, 23 ноември 2001 година) е първият международен
договор в областта~\cite{budapest_convention}. Релевантни са три деяния от дял 1,
посветен на престъпленията срещу поверителността, цялостността и наличността на
компютърните данни и системи.\footnote{Заглавията на членовете на Конвенцията са
дадени в превод на автора, защото официален български текст не беше намерен.
Съдържанието е проверено по английския текст.} Член 2 (незаконен достъп,
\textit{illegal access}) задължава страните да установят като престъпление
умишления достъп до компютърна система или част от нея без право. Страната може да изисква деянието да е извършено
чрез нарушаване на мерки за сигурност или с нечестно намерение. Член 4 (намеса в
данни, \textit{data interference}) обхваща умишленото повреждане, изтриване,
влошаване, променяне или скриване на данни без право. Член 5 (намеса в система,
\textit{system interference}) обхваща умишленото сериозно възпрепятстване на
функционирането на система чрез въвеждане, предаване, повреждане, изтриване,
влошаване, променяне или скриване на данни. За случаи като \textit{Heartland} е
релевантен и член 8 (компютърна измама, \textit{computer-related fraud})~\cite{budapest_convention}.

\subsubsection{Директива 2013/40/ЕС}

Директива 2013/40/ЕС относно атаките срещу информационните системи задължава държавите членки
да предвидят като престъпления сходни деяния и замества Рамково решение 2005/222/ПВР~\cite{eu_2013_40}.
Подредбата на членовете е различна от тази на Конвенцията: в Директивата член 3
урежда незаконния достъп до информационни системи, член 4 незаконната намеса в
система, член 5 незаконната намеса в данни, а член 7 инструментите за извършване на
престъпления. При Конвенцията системата е в член 5, а данните в член 4.
Има и разлика по същество. Член 3 изисква деянието да е наказуемо, когато е
извършено чрез нарушаване на мярка за сигурност, поне в случаите, които не се
считат за маловажни, докато Конвенцията само дава на страните възможност да поставят
това условие. Инжекция, която заобикаля удостоверяването, може да бъде разглеждана
като преодоляване на мярка за сигурност, но това е въпрос на преценка по случая.
Член 9 изисква максимален срок на лишаване от свобода най-малко две години за
деянията по членове 3 до 7 и най-малко пет години за членове 4 и 5, когато са
извършени в рамките на престъпна организация, причиняват сериозни вреди или са
насочени срещу информационна система на критична инфраструктура~\cite{eu_2013_40}.

\subsubsection{Българският Наказателен кодекс}

Компютърните престъпления са уредени в глава девета „а“ „Компютърни
престъпления“, членове 319а до 319е на Наказателния кодекс~\cite{bg_nk}. Текстът е
прочетен на lex.bg към 4 ноември 2026 година и не е сверяван с „Държавен вестник“.
Посочват се само предметът на състава и размерът на лишаването от свобода. Глобите
не се разглеждат.

Член 319а, ал. 1 наказва неправомерния достъп до информационна система или части от
нея, в немаловажни случаи, с лишаване от свобода до шест години. За информация,
представляваща държавна или друга защитена от закон тайна, и за информационна
система, която е част от критична инфраструктура, ал. 4 предвижда от четири до осем
години. Член 319б, ал. 1 наказва неправомерното добавяне, копиране, използване,
променяне, пренасяне, изтриване, повреждане, влошаване, скриване или унищожаване на
компютърни данни или спирането на достъпа до тях, в немаловажни случаи, с лишаване от
свобода до шест години. При цел имотна облага ал. 3 предвижда от пет до седем години.
Член 319в предвижда до седем години, когато деянието по чл. 319б е по отношение на
данни, „които се дават по силата на закон, по електронен път или на електронен или
друг носител“. Член 319г наказва въвеждането на компютърен вирус или на друга
програма, предназначена за нарушаване на дейността на система или за узнаване,
заличаване, изменение или копиране на данни без разрешение (до седем години). Член
319д наказва създаването и разпространението на програми, пароли и кодове за достъп с
цел престъпление по чл. 171, ал. 3 или чл. 319а до 319г (до шест години). Член 319е
се отнася до нарушение на чл. 6, ал. 2, т. 5 от Закона за електронния документ и
електронните удостоверителни услуги~\cite{bg_nk}.

SQL инжекция, използвана за достъп до данни без право, може да осъществи състава на
чл. 319а, ал. 1. Когато чрез нея данни се копират, изменят или изтриват, възможна
квалификация е чл. 319б, ал. 1, а при цел имотна облага и ал. 3. И двата члена
се прилагат в немаловажни случаи. Това е авторов
анализ на текста на закона, а не извод от съдебна практика.

\subsubsection{Администраторът на данни: GDPR}

Българският текст на Регламент (ЕС) 2016/679 (\textit{GDPR}) не съдържа думата
„инжекция“. Член 32, параграф 1
задължава администратора и обработващия да прилагат подходящи технически и
организационни мерки за сигурност, съобразени с риска и с „достиженията на
техническия прогрес“, а буква г) включва процес на редовно изпитване и оценка на
ефективността им. Член 33 изисква уведомяване на надзорния орган без ненужно
забавяне и, когато е осъществимо, не по-късно от 72 часа след узнаването, освен ако
няма вероятност нарушението да породи риск за правата и свободите на лицата. Според
член 83, параграф 4, буква а) нарушенията на задълженията на администратора и
обработващия съгласно членове 8, 11, 25 до 39, 42 и 43, към които спадат членове 32
и 33, подлежат на административно наказание „глоба“ или „имуществена санкция“ в
размер до 10\,000\,000 евро или, в случай на предприятие, до 2\,\% от общия му
годишен световен оборот за предходната финансова година, която от двете суми е
по-висока~\cite{gdpr}. Връзката с инжекциите е тълкувателна: когато за клас уязвимости
съществува документирана и широко известна защита, нейната липса трудно може да се
обоснове с достиженията на техническия прогрес. Така разсъждава \textit{ICO} по
случая \textit{TalkTalk} по действалото тогава право~\cite{ico_talktalk_2016}. Член 33
показва и друго: правото предвижда уведомяване \textit{след} нарушението.

\subsubsection{NIS 2 и Законът за киберсигурност}

Директива (ЕС) 2022/2555 (\textit{NIS 2}) задължава съществените и важните субекти да
приемат подходящи и пропорционални мерки за управление на рисковете за мрежовите и
информационните системи (член 21, параграф 1). Сред минималните мерки е „сигурност
при придобиването на мрежови и информационни системи, разработване и поддръжка“,
включително, в английския текст, „vulnerability handling and disclosure“ (параграф 2,
буква д)~\cite{nis2}. Директивата се прилага към субекти от приложения I и II, които са най-малко
средни предприятия, а в определени случаи независимо от размера (член 2)~\cite{nis2}.
По данни на правна кантора и на проследяващ регистър в България тя е въведена чрез
изменение на Закона за киберсигурност, прието на 5 февруари 2026 година,
обнародвано на 13 февруари 2026 година и в сила от 17 февруари 2026
година~\cite{cms_bg_nis2,risi_bg_nis2}. Заглавната бележка на пречистения текст на
закона потвърждава само изменение, обнародвано в „Държавен вестник“, бр. 17 от 13-о
число на 2026 година~\cite{bg_zks}. Съдържанието на закона по членове не е
проверявано от първичен източник и не се цитира.

\subsection{От наказание към превенция}
\label{sec:cybercrime-prevention}

Разгледаните норми действат след събитието: наказателният закон санкционира
нападателя, регламентът изисква уведомяване след нарушението, а глобата на \textit{ICO}
е последваща. Нито една не намалява вероятността следващото приложение да съдържа
същата слабост.

На 25 март 2024 година \textit{CISA} и \textit{FBI} публикуват съвместно
предупреждение за проектиране с вградена сигурност (\textit{Secure by Design Alert})
„Eliminating SQL Injection Vulnerabilities in Software“. То определя SQL инжекциите
като устойчив клас дефекти, отбелязва, че индустрията знае как да ги елиминира от
десетилетия (като пример: \textit{MySQL} въвежда подготвените изрази през 2004
година), и призовава ръководителите на производителите да прегледат кода си и да
изключат целия клас дефекти от настоящите и бъдещите продукти~\cite{cisa_sqli_alert}.
Основната препоръка е параметризирани заявки с подготвени изрази, които отделят
SQL кода от данните на потребителя. Почистването на входа се оценява като крехко,
трудно за налагане в голям мащаб и често заобиколимо. Препоръчва се
параметризираните заявки да са стандартна практика, наложена в средата за
разработка, например чрез библиотеки, които правят сигурния път пътя по
подразбиране. Документът се позовава и на принципите за проектиране с вградена
сигурност: поемане на отговорност за резултатите за клиентите и изкореняване на
цели класове уязвимости, вместо отстраняване случай по случай~\cite{cisa_sqli_alert}.

Този аргумент е политически, но има техническа последица. Ако сигурният път трябва да
бъде единственият, въпросът е какво може една библиотека да направи невъзможно.
Предупреждението говори за SQL и за стойности на параметри. То не разглежда имената
на таблици и колони, структурата на заявката и хранилищата, които не са релационни.
Именно тук се намира изследователският въпрос на настоящата работа. Следващата глава
разглежда какво вече е предложено като защита и докъде достига то.

\newpage
\section{Преглед на съществуващите решения}
\label{sec:existing}

Предходната глава показа, че инжекциите са добре познат клас уязвимости, чрез който
се извършват престъпления и който води до санкции за администраторите, и че политиката се измества от наказване към изключване на
класа още при проектирането. Настоящата глава разглежда какви средства за това вече
съществуват и къде всяко от тях спира. Първо се разглеждат защитите по слоеве, после
академичната работа върху SQL и инжекциите извън SQL. Следва прегледът на
библиотеките за типово безопасно конструиране на заявки. Главата завършва със
сравнителна таблица и с формулировка на пропуска, който следващите глави адресират.

Прегледът се основава на документацията на проектите и на оригиналните публикации.
Когато в прочетените страници даден факт не е посочен, в таблиците е записано „н.д.“
(не е посочено). Авторът не е стартирал разглежданите библиотеки, затова твърденията
за тях са твърдения на документацията, а не измерване.

\subsection{Защити по слоеве}
\label{sec:existing-layers}

Основният практически източник е наръчникът на \textit{OWASP} за предотвратяване на
SQL инжекции~\cite{owasp_sqli_prevention}. Той поставя на първо място подготвените
изрази с параметризирани заявки, следвани от съхранените процедури и проверката по
списък на допустимите стойности. Екранирането е определено като „STRONGLY
DISCOURAGED“, а към защитите се добавят най-малките привилегии. За всяка защита по-долу
се разглеждат две неща: какво гарантира и как се проваля.

\textbf{Екраниране} е преобразуване на специалните знаци във входа така, че
интерпретаторът да ги третира като данни. Наръчникът настоява ръчното екраниране да не
замества свързването на параметри, защото правилата зависят от диалекта и от
конфигурацията на базата~\cite{owasp_sqli_prevention}. Гаранцията е локална: важи за
конкретния диалект и място, където екранирането е извършено, и се губи на всяко
пропуснато място. Екранирането остава практика там, където параметри не могат да се
ползват: ръководството на \textit{Neo4j} препоръчва за идентификаторите в обратни
апострофи тези знаци да се удвояват~\cite{neo4j_cypher_injection_kb}.

При \textbf{параметризираните заявки} текстът и стойностите се предават поотделно:
заявката съдържа маркери, а драйверът свързва стойностите към тях. Стойността не може
да промени структурата, защото тя вече е разчетена. Това е защитата, препоръчана от
\textit{OWASP}, \textit{CISA} и \textit{FBI}~\cite{owasp_sqli_prevention,cisa_sqli_alert}.
Същият принцип стои зад \textit{Neo4j}, където планът на заявката е компилиран преди
подставянето на стойностите~\cite{neo4j_cypher_injection_kb}. В \textit{Cassandra}
заявка с поне един маркер трябва да бъде подготвена; документацията на \textit{CQL}
описва това като оптимизация и не го свързва със сигурността~\cite{cassandra_cql_prepared}.
Границата е посочена в наръчника на \textit{OWASP}: маркерите заместват стойности, а
не имена на таблици и колони и не ключови думи~\cite{owasp_sqli_prevention}.
Документацията на \textit{EF Core} казва същото~\cite{ef_core_sql}. Защитата не се
прилага автоматично: ако заявката е сглобена с конкатенация, параметри няма. Тя не
отговаря и на въпроса дали потребителят има право да види даден запис.

\textbf{Съхранените процедури} пренасят заявките в базата. Наръчникът предупреждава,
че свързването на аргументите защитава извикването, но не и динамичния SQL, който
процедурата сглобява вътре в себе си~\cite{owasp_sqli_prevention}. Халфонд, Виегас и
Орсо описват същото като разпространено заблуждение~\cite{halfond2006classification}.
Гаранцията е условна: зависи от начина, по който е написана процедурата.

Проверката по \textbf{списък на допустимите стойности} (\textit{allow-list}) се
препоръчва за структурни избори, които не могат да се параметризират. Пример е
посоката на сортиране: входът се съпоставя с литералите \code{ASC} или \code{DESC},
притежавани от приложението. Проверката допълва свързването на параметри, но не го
замества~\cite{owasp_sqli_prevention}. Гаранцията е силна само когато множеството може
да се изброи и не помага при свободен текст.

\textbf{Най-малките привилегии} ограничават какво може да направи вече променена
заявка. \textit{OWASP} предупреждава приложението да не получава права на
администратор на базата~\cite{owasp_sqli_prevention}, а \textit{Neo4j} препоръчва
контрол на достъпа по роли~\cite{neo4j_cypher_injection_kb}. Мярката не предотвратява
инжекцията, а намалява последиците ѝ.

\textbf{Защитната стена за уеб приложения} (\textit{web application firewall},
\textit{WAF}) проверява HTTP трафика пред приложението. \textit{ModSecurity} е
отворен механизъм за такива стени за \textit{Apache}, \textit{IIS} и
\textit{Nginx}~\cite{modsecurity_repo}, а \textit{OWASP Core Rule Set} (\textit{CRS})
е наборът правила за него. Всяко съвпаднало правило добавя точки към резултата на
заявката, която се отхвърля при праг, по подразбиране 5~\cite{crs_anomaly}.
По-високите от четирите нива на параноя откриват повече атаки, но дават повече
фалшиви тревоги, които налагат изключения на правила и могат да доведат до
изключване на \textit{CRS}~\cite{crs_paranoia}. Стената вижда HTTP заявката, а не
заявката към базата, която приложението по-късно ще сглоби. Гаранцията ѝ е
вероятностна и не променя причината за уязвимостта.

Защитите се разделят на две групи. Екранирането, параметрите и списъкът на допустимите
стойности действат върху самата заявка и дават силна гаранция, но само за онази част,
за която са приложими, и само когато програмистът ги приложи. Привилегиите и
защитната стена действат извън заявката и ограничават вредата или вероятността.
Общото между всички е, че се прилагат от човек на всяко място поотделно.

\subsection{Академична работа по откриване и предотвратяване}
\label{sec:existing-academic}

Су и Васерман определят инжекцията формално: за да успее атака, входът, който попада
в заявката, трябва да променя нейната предвидена синтактична структура. Според тях
това е първото формално определение на атаките чрез инжектиране на команди в уеб
приложенията. Заедно с него авторите дават алгоритъм за предотвратяването им, основан
на контекстно-свободни граматики. Реализацията им \textit{SQLCHECK}, оценена върху
реални приложения със системно събрани реални атаки, по думите им не е дала нито
фалшиво положителни, нито фалшиво отрицателни резултати~\cite{su2006essence}. Определението
се използва и в настоящата работа (глава~\ref{sec:theory}).

Халфонд, Виегас и Орсо правят обзор на известните видове SQL инжекции (сред тях
тавтологии, обединени и натрупани заявки, атаки чрез съхранени процедури,
алтернативни кодировки) и оценяват съществуващите техники спрямо тях. Основното им
наблюдение е, че много предложени решения адресират само подмножество от
проблема~\cite{halfond2006classification}. Техниката \textit{AMNESIA} на Халфонд и
Орсо строи чрез статичен анализ модел на законните заявки на приложението и проверява
всяка генерирана заявка спрямо него по време на изпълнение. Върху седем приложения
тя е спряла всички опитани атаки без фалшиво положителни резултати~\cite{halfond2005amnesia}.
\textit{CANDID} на Бандхакави и съавтори е, според резюмето, схема за
автоматично преобразуване на уеб приложения, така че по твърдение на авторите те да
са безопасни срещу SQL инжекции. Методът е посочен в заглавието като динамична
оценка на кандидати~\cite{bandhakavi2007candid}. \textit{SQLrand} на Бойд и
Керомитис рандомизира набора от инструкции на SQL, така че инжектираният текст да е
невалиден за парсера, а междинен прокси превежда рандомизирания език към стандартния
за \textit{MySQL}~\cite{boyd2004sqlrand}.

Тези работи имат общи черти. Всички са насочени към SQL. Действат върху вече
построена заявка: сравняват я с модел, с кандидат или с рандомизиран речник, а не
променят начина, по който се построява. Изискват допълнителен инструмент: анализ,
модифицирана среда или прокси. Гаранцията се налага по време на изпълнение върху вече
построената заявка, а не при построяването ѝ. Някои от работите дават и формално
обосновани алгоритми (така Су и Васерман претендират за алгоритъм, който е коректен и
пълен), но оценката им е върху подбрани приложения. Определението на Су и Васерман
остава общата основа, спрямо която се измерва успехът на всички тях.

\subsection{Инжекции извън SQL}
\label{sec:existing-nosql}

Литературата опровергава мнението, че хранилищата без SQL са имунизирани. Рон,
Шулман-Пелег и Бронщайн представят техники за инжектиране в NoSQL хранилища, сред
които инжектиране чрез масиви в \textit{PHP}, \textit{OR} инжектиране в
\textit{MongoDB} и инжектиране на произволен \textit{JavaScript}. В един пример с
оператора \code{$ne} условията за потребителско име и парола се изпълняват за всяка
стойност, различна от посочената~\cite{ron2015nosql}. Част от същите автори публикуват
статия в \textit{IEEE Security \& Privacy}, според резюмето на която само анализът на
кода е недостатъчен за предотвратяване на такива атаки и мерките трябва да обхващат
целия жизнен цикъл на софтуера~\cite{ron2016nosql}. Ръководството за тестване на уеб сигурността на
\textit{OWASP} (версия 4.2) обяснява, че атаките срещу NoSQL могат да се изпълняват в
процедурен език, затова въздействието им може да е по-голямо, и че филтрирането на
знаци, специални за HTML, не защитава JSON интерфейс. В \textit{MongoDB} интерфейсът
очаква BSON и предоставя безопасно сглобяване на заявки, но несериализираният JSON и
изразите на \textit{JavaScript} се допускат в алтернативни параметри, като най-често
използваният е операторът \code{$where}~\cite{owasp_wstg42_nosql}.

Документацията на \textit{Neo4j} назовава Cypher инжекцията „братовчед“ на SQL
инжекцията и посочва като основна защита подаването на входа като параметри.
Конкатенирането на вход в низ, подаван на \code{apoc.cypher.doIt}, остава уязвимо,
а етикетите, типовете на връзките и имената на свойства според статията не могат да
се подават като параметри и затова се препоръчва да се избягват или санитизират~\cite{neo4j_cypher_injection_kb}.
Текущото ръководство по Cypher уточнява, че параметри не могат да заместват директно
ключове на свойства, типове връзки и етикети (например \code{n:$param} е невалидно),
но етикетите и типовете могат да се задават динамично чрез израз~\cite{neo4j_cypher_parameters}.
Следователно в Cypher имената остават избираеми по време на изпълнение по документиран
път.

Документацията на \textit{Redis} приема, че сървърът е за доверени клиенти в доверена
среда и че недоверен достъп трябва да се опосредства от слой, проверяващ входа. За
инжекциите тя твърди, че протоколът няма понятие за екраниране на низове, че „injection
is impossible under normal circumstances using a normal client library“ и че протоколът
използва низове с префикс за дължина, безопасни за двоични данни. Същото важи за
скриптовете на \textit{Lua}, но тялото на скрипт не бива да се сглобява от недоверени
низове~\cite{redis_security}. Условието „нормална клиентска библиотека“ означава, че
гаранцията на протокола не е гаранция на приложението. За \textit{Cassandra}
документацията на CQL описва маркерите и подготвените изрази, но в прочетените
страници не е открито самостоятелно ръководство за инжекции в CQL~\cite{cassandra_cql_prepared}.

Всяко хранилище решава проблема със свои средства: параметри, подготвени изрази,
протокол с префикс за дължина. Стойностите са защитени по замисъл на протокола или
интерфейса, но точно там, където приложението избира структура или име,
документацията предупреждава за граница. Гаранциите са формулирани поотделно от
производителя на всяко хранилище. В прегледаните източници няма общо определение,
което да обхваща всички хранилища, а академичните работи от предходния раздел не
покриват нито едно от тях.

\subsection{Типово безопасно конструиране на заявки}
\label{sec:existing-typesafe}

Друг подход е заявките да не се пишат като низове, а да се изграждат от типове и
функции на езика. Така компилаторът проверява имената и типовете. Разглеждат се
\textit{LINQ} с \textit{EF Core} (C\#), \textit{jOOQ} (Java), \textit{Diesel} и
\textit{sqlx} (Rust), \textit{esqueleto} и \textit{opaleye} (Haskell) и
\textit{sqlpp11} (C++). За всяка се търси какво прави безопасно и дали остава път към
суров текст или към идентификатори, избрани по време на изпълнение.

Според документацията на \textit{EF Core} методите \code{FromSql} и
\code{FromSqlInterpolated} са безопасни, защото предават данните като параметър, а
\code{FromSqlRaw} може да е уязвим. Когато името на колона трябва да се вмъкне
динамично, се ползва \code{FromSqlRaw}, а санитизирането е отговорност на
програмиста~\cite{ef_core_sql}. В \textit{jOOQ} SQL обикновено се строи като типизирано дърво със
свързани стойности и според документацията по този път не е възможно да се въведат инжекции. Изключение
са методите за суров SQL, отбелязани с \code{org.jooq.PlainSQL} от версия 3.6~\cite{jooq_sqli},
чието използване документацията свързва с проверка чрез \textit{Checker Framework} или
\textit{Error Prone}~\cite{jooq_plain_sql}. \textit{Diesel} заявява стремеж да
елиминира неправилните взаимодействия с базата на етапа на компилацията~\cite{diesel_site}
и според документацията изисква схемата да е описана в кода~\cite{diesel_crate}. Суров
SQL се пише чрез \code{sql\_query}, а малки фрагменти чрез \code{sql}; страницата на
\code{sql\_query} не споменава инжекциите и предупреждава само за типовете~\cite{diesel_sql_query}.
Макросът \code{query!} на \textit{sqlx} проверява SQL низа при компилация срещу жива
база или запазен кеш~\cite{sqlx_readme}. Във версия 0.9.0 функцията \code{sqlx::query}
приема \code{impl SqlSafeStr}, динамично построен низ трябва изрично да се обвие в
\code{AssertSqlSafe}, а документацията предупреждава срещу \code{format!}~\cite{sqlx_query_docs}.

В \textit{esqueleto} има \code{rawSql} с текст и списък от стойности, а SQL функция
може да се въведе чрез \code{unsafeSqlFunction} или \code{unsafeSqlBinOp}~\cite{esqueleto_hackage}.
В \textit{opaleye} функциите за таблица и колона приемат обикновен низ:
\code{table :: String -> ...} и \code{tableField :: String -> ...}~\cite{opaleye_table},
а модулът съдържа и \code{unsafeCast :: String -> ...}~\cite{opaleye_main}. Това не
доказва, че недоверени имена могат да бъдат подадени, но показва, че типовата система
не го забранява. \textit{sqlpp11} е типово безопасен вграден език за SQL в C++, чиито
таблици се пораждат от DDL и чиито заявки се проверяват при компилация за грешки в
синтаксиса, типовете и имената~\cite{sqlpp11_repo}. За неподдържани изрази
документацията описва \code{sqlpp::verbatim}~\cite{sqlpp11_functions}. Същият автор
поддържа и \textit{sqlpp23}, чиято документация описва \code{verbatim_table}, който
добавя таблици, дори когато те не са известни при компилация, както и \code{verbatim} и
\code{verbatim_clause}~\cite{sqlpp23_custom,sqlpp23_repo}. Това е най-близката
прегледана библиотека на C++ до подхода на настоящата работа и затова е важно, че и в
нея път към суров текст е предвиден по дизайн. Преглед, изчерпателен за всички
библиотеки на C++, не е правен.

\begin{table}[!htbp]
\centering
\caption{Типово безопасни библиотеки: какво защитават и къде остава изход}
\label{tab:libs}
\footnotesize
\setlength{\tabcolsep}{3pt}
\begin{tabularx}{\textwidth}{L{1.9cm}L{2.6cm}L{3.2cm}L{3.5cm}Y}
\toprule
\textbf{Библиотека} & \textbf{Стойности} & \textbf{Идентификатори и структура} & \textbf{Път към суров текст} & \textbf{Проверка} \\
\midrule
EF Core (LINQ) & Параметри при \code{FromSql}~\cite{ef_core_sql} & Имена на колони не се параметризират, нужен е \code{FromSqlRaw}~\cite{ef_core_sql} & \code{FromSqlRaw}, \code{SqlQueryRaw}, \code{ExecuteSqlRaw} & Типове при компилация, превод в SQL при изпълнение \\
jOOQ & Свързани стойности в дървото~\cite{jooq_sqli} & Типизирано дърво~\cite{jooq_sqli} & Методи с \code{@PlainSQL}~\cite{jooq_sqli,jooq_plain_sql} & Компилация (Java), по желание допълнителна проверка \\
Diesel & Свързани стойности~\cite{diesel_debug_query} & Схемата е описана в кода~\cite{diesel_crate} & \code{sql\_query}, \code{sql} & Компилация \\
sqlx & Свързани параметри~\cite{sqlx_query_docs} & \code{query!} проверява литерален низ срещу база~\cite{sqlx_readme} & \code{AssertSqlSafe} върху динамичен низ~\cite{sqlx_query_docs} & Компилация срещу жива база или кеш \\
esqueleto & \code{rawSql} приема стойности, за DSL н.д.~\cite{esqueleto_hackage} & н.д. & \code{rawSql}, \code{unsafeSqlFunction}, \code{unsafeSqlBinOp}~\cite{esqueleto_hackage} & Компилация (Haskell) \\
opaleye & н.д. & Имена като \code{String}~\cite{opaleye_table} & \code{unsafeCast} с низ~\cite{opaleye_main} & Компилация (Haskell) \\
sqlpp11 / sqlpp23 & В sqlpp11 подготвени изрази с \code{parameter()}~\cite{sqlpp11_prepared} & В sqlpp11 таблици от DDL~\cite{sqlpp11_repo}, в sqlpp23 е възможен \code{verbatim_table} & \code{verbatim}, \code{verbatim_table}, \code{verbatim_clause}~\cite{sqlpp11_functions,sqlpp23_custom} & Компилация (C++) \\
\bottomrule
\end{tabularx}
\end{table}

Таблица~\ref{tab:libs} показва повтарящ се модел. Защитата на стойностите е нормален
случай там, където документацията я посочва. Идентификаторите и структурата са
защитени дотолкова, доколкото програмистът остава в типизирания интерфейс. Във всяка
прегледана библиотека обаче е документиран път към суров текст. Обозначен е различно:
наставката \code{Raw} в \textit{EF Core}, анотацията \code{@PlainSQL} с възможност за
проверка в \textit{jOOQ}, явното \code{AssertSqlSafe} в \textit{sqlx}. Тези
обозначения помагат при преглед, но налагането им е конвенция или външен инструмент, а
не свойство на типовата система. В \textit{opaleye} именуването чрез низове е част от
обикновения интерфейс. Таблицата не казва, че библиотеките са небезопасни, а че
гаранцията им зависи от това, че изходът няма да бъде използван.

\subsection{Сравнение на подходите}
\label{sec:existing-comparison}

Таблица~\ref{tab:comparison} съпоставя разгледаните подходи по критерии, които
следват от определението на Су и Васерман~\cite{su2006essence}: кои части на заявката
се защитават, покриват ли се хранилища извън SQL, кога се налага гаранцията и има ли
изход. Оценките са на автора, въз основа на цитираните по-горе източници.
„Частично“ означава, че гаранцията важи при определени условия, посочени в
съответния раздел.

\begin{table}[!htbp]
\centering
\caption{Сравнение на подходите по критерии}
\label{tab:comparison}
\footnotesize
\setlength{\tabcolsep}{3pt}
\begin{tabularx}{\textwidth}{L{2.4cm}L{2.0cm}L{2.5cm}L{2.0cm}L{1.9cm}L{1.9cm}Y}
\toprule
\textbf{Подход} & \textbf{Стойности} & \textbf{Иденти-\newline фикатори} & \textbf{Структура} & \textbf{Не-SQL хранилища} & \textbf{Кога се налага} & \textbf{Изход} \\
\midrule
Екраниране & частично & частично & не & зависи от диалекта & изпълнение & всяко пропуснато място \\
Подготвени изрази & да & не & не & да, където има маркери & изпълнение & сглобяване на низ \\
Съхранени процедури & при параметри вътре & не & не & н.д. & изпълнение & динамичен SQL вътре \\
Списък на допустимите & да, за изброими & да, за изброими & да, за изброими & да & изпълнение & свободен текст \\
Най-малки привилегии & ограничава вредата & ограничава вредата & ограничава вредата & да, като принцип & внедряване & не е приложимо \\
Защитна стена & по вероятност & по вероятност & по вероятност & според правилата & внедряване & пропуснати атаки \\
Откриване по структура (SQLCHECK, AMNESIA, CANDID, SQLrand) & да & да & да & не, само SQL & изпълнение, с инструмент & само SQL в оценките \\
Типови езици за заявки & да & зависи от библиотеката & да, без суров текст & основно SQL & компилация & да, във всички прегледани \\
\bottomrule
\end{tabularx}
\end{table}

Таблица~\ref{tab:comparison} позволява три наблюдения. Първо, никой подход не
съчетава защита на стойностите, идентификаторите и структурата, покритие на хранилища
извън SQL и липса на изход. Параметрите покриват стойностите, но по природа не покриват
имената. Откриването по структура покрива всички части на SQL заявката, но е насочено
само към SQL и действа върху вече построена заявка. Типовите езици за заявки са
единствените, които налагат гаранцията при компилация, но във всички прегледани
случаи оставят път към суров текст. Второ, подходите, които действат при изпълнение
или внедряване, откриват или блокират проблема, след като кодът вече съществува, а
подходите при компилация могат да предотвратят самото изграждане на небезопасната
заявка, но гаранцията им е толкова силна, колкото е затворен изходът. Трето, ако
изходът е част от гаранцията, нито един прегледан подход не дава гаранция за целия
клас, а само за частта, която съответства на неговия механизъм. Тази разлика между
защита на отделните случаи и изключване на класа е същността на аргумента на
\textit{CISA} и \textit{FBI} (раздел~\ref{sec:cybercrime-prevention}).

\subsection{Пропускът}
\label{sec:existing-gap}

В рамките на прегледаната литература и документация могат да се направят три
наблюдения. Първо, работата по откриване и предотвратяване на инжекции се
концентрира върху SQL~\cite{su2006essence,halfond2006classification,halfond2005amnesia,bandhakavi2007candid,boyd2004sqlrand}.
Работата върху NoSQL е по-малко и описва техники и мерки за смекчаване~\cite{ron2015nosql,ron2016nosql},
а документацията за Neo4j, Redis и Cassandra е формулирана отделно за всяко
хранилище. Второ, типовите езици за заявки защитават стойностите, но в прегледаните
библиотеки запазват начин да се подаде идентификатор или текст на заявката по време на
изпълнение. Той е обозначен по различен начин, но е налице. Трето, \textit{SQLCHECK} и
\textit{AMNESIA} са измерени върху реални приложения, но само за SQL~\cite{su2006essence,halfond2005amnesia}.
В прегледаните източници не е открита оценка, която прилага гаранция, наложена при
компилация, към няколко разнородни хранилища и я сравнява с измерена уязвима базова
реализация.

Пропускът е следователно в съчетанието на три условия, а не във всяко от тях
поотделно: гаранция, наложена при компилация, обхващаща няколко разнородни хранилища
и оценена спрямо измерена уязвима базова реализация. Формулировката е ограничена до
прегледаната литература и не е твърдение за абсолютна новост. Следващите глави
адресират този пропуск: глава~\ref{sec:theory} въвежда понятията, а останалите
разглеждат подхода, неговата реализация и оценката му.

\newpage
\section{Теоретична основа}
\label{sec:theory}

Предходната глава позиционира работата спрямо съществуващите защити. Настоящата
глава въвежда само понятията, необходими за разбирането на собствения принос:
формално определение на инжекцията, трите принципно различни начина за нейното
предотвратяване, разликите между езиците за заявки на изследваните хранилища и
ролята на времето на компилация като граница на доверие. Последната точка е
мостът между определението и реализацията, разгледана в
глава~\ref{sec:enforcement}.

\subsection{Инжекцията като промяна на синтактичната структура}
\label{sec:theory-definition}

Интуитивното описание на инжекцията („въведеният текст се изпълнява като код“) е
недостатъчно точно, за да служи за критерий при оценка. В настоящата работа се
използва определението на Су и Васерман~\cite{su2006essence}. Заявката се
разглежда като изречение от формален език и съответно като синтактично дърво.
Всеки фрагмент от заявката, произлизащ от недоверен вход, трябва да попада изцяло
в едно листо на това дърво, тоест да бъде стойност. \textbf{Инжекция} е налице
тогава, когато недоверен вход променя синтактичната структура на заявката: когато
той обхваща повече от едно листо или въвежда нови вътрешни възли.

Определението има две предимства. Първо, то е независимо от конкретния полезен
товар (\textit{payload}): то не изисква списък от опасни символи, а критерий за
структурна цялост. Второ, то е независимо от езика. Едно и също понятие обхваща
SQL, документните заявки на MongoDB, Cypher, CQL и командния протокол на Redis,
стига за всеки от тях да е известна граматиката, по която хранилището разбира
заявката. Именно тази независимост позволява изследователският въпрос да бъде
поставен за цяло семейство от класове инжекции, а не за всеки поотделно.

Класификацията на Халфонд, Виегас и Орсо~\cite{halfond2006classification}
допълва определението от страна на практиката. Тя разграничава атаките по
механизъм (тавтологии, обединения чрез \code{UNION}, натрупани заявки, заявки,
предизвикващи грешка, слепи атаки по истинност и по време) и по намерение
(извличане на данни, заобикаляне на удостоверяване, промяна на данни, отказ от
услуга). Корпусът от полезни товари, използван в
глава~\ref{sec:evaluation}, е подреден по същите механизми.

\subsection{Три начина за отделяне на данните от кода}
\label{sec:theory-remedies}

Общата причина за всички инжекции е една: данни и код се пренасят по един и същ
канал, най-често като един низ, и хранилището ги разделя едва при синтактичния
анализ. Всяка защита следователно отговаря на въпроса как недоверената стойност
да остане листо в дървото. Принципно се различават три подхода.

\textbf{Екраниране} (\textit{escaping}). Стойността остава в същия канал, но се
преобразува така, че синтактичният анализатор да не я разпознае като граница на
листото. Сигурността зависи от пълното и правилно преобразуване за конкретния
диалект и кодиране на символите. Пропуск на един специален символ или разлика
между представата на приложението и на хранилището за кодирането е достатъчна,
за да се загуби гаранцията.

\textbf{Свързване на параметри} (\textit{parameter binding}). Заявката се изпраща
с маркери на местата на стойностите, а самите стойности пътуват отделно, по
протокол, който хранилището никога не подлага на синтактичен анализ като заявка.
Стойността не може да стане вътрешен възел, защото не участва в анализа изобщо.
Това е механизмът на подготвените заявки~\cite{owasp_sqli_prevention}.

\textbf{Конструиране} (\textit{construction}). Структурата на заявката не се
сглобява от низове, а се изразява в средство, в което недоверените данни нямат
път към структурата. Свързването защитава стойностите, но оставя отворен въпроса
кой определя структурата. Конструирането отговаря на този въпрос.

Разграничението е съществено, защото свързването не покрива всички части на
заявката. Нито едно от изследваните хранилища не позволява параметър на мястото
на име на таблица или колона. Затова приложения, които използват подготвени
заявки за стойностите, нерядко сглобяват идентификаторите като низове, и именно
там остава уязвимостта. Настоящата работа изследва конструирането като допълнение
към свързването, което затваря тази празнина.

\subsection{Езиците на изследваните хранилища}
\label{sec:theory-languages}

Определението от точка~\ref{sec:theory-definition} предполага граматика. Петте
класа инжекции в обхвата на работата се различават именно по това каква
граматика анализира хранилището и дали изобщо анализира текст, получен от
приложението.

\begin{table}[H]
\centering
\caption{Граматика и клас инжекция при изследваните хранилища}
\label{tab:theory-languages}
\begin{tabularx}{\textwidth}{L{2.6cm}L{3.2cm}Y}
\toprule
\textbf{Хранилище} & \textbf{Какво се анализира} & \textbf{Клас инжекция} \\
\midrule
PostgreSQL, MySQL, SQLite & текст на SQL & SQL инжекция: тавтологии, \code{UNION}, натрупани заявки, слепи атаки \\
Cassandra & текст на CQL & CQL инжекция, сходна със SQL, но с по-беден език \\
Neo4j & текст на Cypher & Cypher инжекция, включително извикване на процедури \\
MongoDB & документ BSON, не текст & инжекция на оператори: стойност, която се интерпретира като {\ttfamily \$ne} или {\ttfamily \$where} \\
Redis & кадри RESP с префикс за дължина & инжекция на протоколно ниво чрез \code{CR LF}, ако командата се сглобява като текст \\
\bottomrule
\end{tabularx}
\end{table}

Таблица~\ref{tab:theory-languages} показва, че хранилищата не са еднородни по
отношение на заплахата. При SQL, CQL и Cypher хранилището анализира текст и
опасността е класическата. При MongoDB заявката е структуриран документ и
инжекцията се изразява по друг начин: стойност, която приложението очаква да е
низ, пристига като обект с ключ, започващ с {\ttfamily \$}, и се интерпретира като
оператор. При Redis протоколът RESP предава всеки аргумент с изрично указана
дължина. Ако клиентът кодира аргументите правилно, съдържанието на аргумента не
може да бъде прочетено като граница на команда, независимо от символите в него.
Следователно за две от седемте хранилища правилното кодиране премахва цялата
граматика, от която атакуващият би могъл да излезе, а не само я екранира.

\subsection{Времето на компилация като граница на доверие}
\label{sec:theory-compile-time}

Конструирането от точка~\ref{sec:theory-remedies} изисква средство, в което
недоверените данни нямат път към структурата. В езика C++ такова средство е
разграничението между стойностите, известни по време на компилация, и тези,
известни едва по време на изпълнение.

Аргументите на шаблон, които не са типове (\textit{non-type template
parameters}), трябва да бъдат \textbf{константни изрази}. Стойност, която
зависи от входа на програмата, по дефиниция не е константен израз и не може да
бъде такъв аргумент. Ако следователно всеки елемент на синтактичната структура
на заявката (име на таблица, име на колона, вид на операцията, оператор на
сравнение, наличие на клауза) е закодиран като аргумент на шаблон или като тип,
то никоя стойност, получена по време на изпълнение, не може да го промени. Входът
може да достигне заявката само там, където типовата система изрично го допуска.

Аргументът може да се формулира като свързване на определението с типовата
система:

\begin{enumerate}
    \item Нека всички вътрешни възли на синтактичното дърво на заявката са
          определени от аргументи на шаблони и типове. Тогава те са фиксирани в
          момента на компилация.
    \item Нека единственият път, по който стойност по време на изпълнение достига
          заявката, е типизиран заместител, който се свързва като параметър
          съгласно точка~\ref{sec:theory-remedies}.
    \item Тогава недоверен вход може да заеме само листо и не може да промени
          структурата. Съгласно определението на Су и
          Васерман~\cite{su2006essence} инжекция не е възможна.
\end{enumerate}

Двете предпоставки не се приемат на доверие. Първата се проверява в
глава~\ref{sec:evaluation} чрез програми, които се опитват да подадат стойност по
време на изпълнение на мястото на идентификатор или структура и трябва да бъдат
отхвърлени от компилатора. Втората се проверява чрез диференциален тест и чрез
анализа на всеки конектор в глава~\ref{sec:enforcement}. Именно проверката на
втората предпоставка разкри случай, в който тя не е била изпълнена: казусът в
глава~\ref{sec:case-study}.

Цената на подхода е симетрична на гаранцията. Ако идентификаторите и структурата
не могат да зависят от входа, то и програмистът не може да ги направи зависими.
Схеми, които се определят динамично, не се поддържат. Това не е пропуск, а
същността на гаранцията, и е обсъдено като ограничение в
глави~\ref{sec:threat-model} и в заключението.

\newpage
\section{Модел на заплахите и метод}
\label{sec:threat-model}

Теоретичната основа от глава~\ref{sec:theory} показа при какви предпоставки
конструирането по време на компилация изключва инжекцията. За да бъде това
твърдение проверимо, е необходимо да се определи срещу какъв нападател се
твърди гаранцията, какво точно се защитава и по какви критерии ще се прецени
успехът. Настоящата глава фиксира тези три неща преди представянето на каквито и
да било резултати и описва метода, по който те се проверяват.

\subsection{Защитаван актив}
\label{sec:tm-asset}

Защитаваният актив е \textbf{структурата на командата}, която приложението
изпраща към хранилището за данни. Стойностите в една заявка са данни, а
структурата определя какво хранилището е упълномощено да върне или промени.
Нападател, който превръща {\ttfamily WHERE id = \$1} в
{\ttfamily WHERE id = \$1 OR 1=1}, не е променил данни. Той е променил обхвата на
правомощията на заявката. Всеки клас инжекция от
таблица~\ref{tab:theory-languages} е разновидност на този ход.

\subsection{Модел на нападателя}
\label{sec:tm-attacker}

Нападателят контролира вход, който приложението използва при изграждането на
заявка. Той не контролира изходния код на приложението, компилатора и
драйвера на хранилището. Разграничават се три нива по нарастваща сила.
Разграничението е съществено, защото обичайното обсъждане на инжекциите ги
слива.

\textbf{Ниво 1: нападателят контролира стойности.} Стойност от формуляр, заглавна
част на заявка или съобщение достига заявката. Това е обичайният случай и
единственият, който свързването на параметри покрива.

\textbf{Ниво 2: нападателят контролира идентификатори.} Нападателят избира име на
таблица, на колона или ключ за подреждане. Тук свързването на параметри не
помага, защото нито едно от изследваните хранилища не позволява параметър на
мястото на идентификатор.

\textbf{Ниво 3: нападателят контролира структурата.} Нападателят определя формата
на заявката, например кои условия участват в клаузата \code{WHERE}, чрез
построител на филтри, който сглобява предикати според заявката.

Гаранцията на изследваната библиотека се твърди и за трите нива. На ниво 1 тя
разчита на свързване, а на нива 2 и 3 на това, че съответната възможност
изобщо не е изразима по време на изпълнение (точка~\ref{sec:theory-compile-time}).
Изключение са миграциите, където защитата на ниво 2 е проверка по време на
изпълнение (точка~\ref{sec:enf-ddl}).

\subsection{Извън обхвата}
\label{sec:tm-scope}

Следните заплахи съзнателно са извън твърдяната гаранция и се обсъждат като
ограничения, а не като пропуски:

\begin{itemize}
    \item \textbf{инжекция от втори ред в приложението}: стойност, прочетена от
          хранилището и използвана от приложението като идентификатор другаде, е
          извън видимостта на библиотеката. Собственият случай на библиотеката,
          миграция, която чете имена обратно от хранилището, е в обхвата
          (точка~\ref{sec:enf-ddl});
    \item \textbf{логически грешки в приложението}: липсваща проверка на
          правомощия или прекалено широк филтър не са инжекция;
    \item \textbf{уязвимости в драйверите} (\code{libpq}, \code{hiredis},
          \code{bsoncxx} и останалите), които се приемат за коректни;
    \item \textbf{отказ от услуга} чрез скъпи, но легитимни заявки, който е
          въпрос на капацитет, а не на инжекция.
\end{itemize}

\subsection{Изследователски въпрос и хипотеза}
\label{sec:tm-question}

Работата отговаря на въпроса:

\begin{quote}
\textit{Може ли една абстракция, конструираща заявките по време на компилация, да
елиминира цяло семейство от класове инжекции при хетерогенни хранилища за
данни, и къде точно спира тази гаранция?}
\end{quote}

Хипотезата е, че при изпълнени предпоставки от
точка~\ref{sec:theory-compile-time} никой полезен товар, подаден през
документирания интерфейс, няма наблюдаем ефект върху хранилището, докато същият
полезен товар, подаден в наивно сглобена заявка, има. Втората част на въпроса,
къде спира гаранцията, е равностойна на първата. Работа, която отчита само
успехи, не позволява да се прецени валидността на твърденията си.

\subsection{Критерии за успех}
\label{sec:tm-criteria}

Критериите са определени преди измерванията и не са променяни след тях.

\begin{table}[H]
\centering
\caption{Критерии за успех}
\label{tab:tm-criteria}
\begin{tabularx}{\textwidth}{L{1.0cm}Y}
\toprule
\textbf{№} & \textbf{Критерий} \\
\midrule
К1 & Никой полезен товар от корпуса, подаден през документирания интерфейс, не
     предизвиква наблюдаем ефект: върнати редове, които не би трябвало да
     съществуват, промяна в данните или схемата, или принудително изпълнение. \\
К2 & Всеки опит идентификатор или структура на заявката да зависи от стойност по
     време на изпълнение се отхвърля от компилатора, и то по причината, която се
     проверява, а не по случайна друга. \\
К3 & Наивната заявка, сглобена от същите полезни товари, е компрометирана поне
     веднъж при същата постановка. Без това изпълнението на К1 не доказва нищо. \\
К4 & Границите на гаранцията са назовани и, където е възможно, демонстрирани. \\
\bottomrule
\end{tabularx}
\end{table}

Критерий К3 заслужава коментар. Сравнение, при което нито защитеният, нито
незащитеният път са засегнати, не различава защитата от неефективен корпус.
Наивната заявка е съзнателно уязвима и служи за база, спрямо която резултатът на
библиотеката придобива смисъл. Ако тя не бъде компрометирана, тестът се отчита
като неуспешен.

\subsection{Метод на оценката}
\label{sec:tm-method}

Методът съчетава три независими източника на доказателства, всеки от които
покрива различна предпоставка.

\textbf{Отрицателни компилационни тестове} (за К2). Всеки тест е програма, която
се опитва да направи идентификатор или структура зависими от стойност по време на
изпълнение и трябва да бъде отхвърлена от компилатора. Такъв тест сам по себе си
е слаб: той би преминал и при грешка, несвързана с проверяваното свойство,
например липсващ заглавен файл. Затова се прилагат две мерки. Първо, всяко
отхвърляне трябва да съответства на образец за очакваната диагностика. Второ,
наборът съдържа \textbf{контролни случаи}, които използват същата постановка и
трябва да се компилират. Ако постановката е повредена, контролните случаи също
се провалят.

\textbf{Диференциално тестване} (за К1 и К3). Всеки приложим полезен товар се
подава по два пътя, които задават един и същ въпрос към една и съща база данни:
през наивно сглобена заявка и през библиотеката. За всяко изпълнение се използва
нова база данни. Резултатът се определя от \textbf{оракул}, който наблюдава само
ефекти: върнати редове, промяна на състоянието и принудително прекратяване на
изпълнението.

\textbf{Анализ на изходния код} (за предпоставка 2 от
точка~\ref{sec:theory-compile-time} и за К4). За всеки от седемте конектора се
установява как точно стойността достига драйвера.

\subsection{Какво методът не доказва}
\label{sec:tm-not-proven}

Методът има граници, които се посочват предварително, за да не се надценяват
резултатите.

\begin{itemize}
    \item Оракулът вижда само ефекти. Сляпа атака, която променя логиката на
          заявката, но връща нула редове, не се отчита като компрометиране.
          Следователно броят на компрометиранията на наивния път е \textbf{долна
          граница}. Това изместване работи срещу твърдението на работата, а не в
          негова полза.
    \item Диференциалният тест използва една форма на заявка за хранилище. Друга форма би
          могла да позволи други атаки срещу наивния път.
    \item Корпусът е краен. Отсъствие на успешна атака срещу библиотеката не е
          доказателство за отсъствие на всякаква атака. Доказателството за това се
          търси в аргумента от точка~\ref{sec:theory-compile-time} и в
          компилационните тестове, а не в корпуса.
    \item Компилационните тестове се изпълняват с GCC и Clang. Поведението на
          други компилатори не е проверено.
\end{itemize}

\newpage
\section{Реализация на гаранцията при седем хранилища}
\label{sec:enforcement}

Моделът на заплахите от глава~\ref{sec:threat-model} изисква гаранцията да се
провери поотделно за всяко хранилище, защото предпоставка 2 от
точка~\ref{sec:theory-compile-time} се изпълнява от конкретния конектор, а не от
общата абстракция. Настоящата глава описва накратко платформата, установява
откъде идва всяка част на заявката и проследява за всеки от седемте конектора как
стойността достига драйвера. Номерата на редовете се отнасят към версия
\code{a17f4e7} на хранилището с изходния код.

\subsection{Платформата}
\label{sec:enf-platform}

Изследваната библиотека е разработена в предходна дипломна работа за ОКС
„Бакалавър“~\cite{tsvetanov2026orm}, в която са описани обектният модел,
междинното представяне на заявките, асинхронното изпълнение и конекторите.
Настоящата работа не преповтаря тези резултати и не претендира за тях. Тя
използва библиотеката като платформа и изследва въпрос, който предходната работа
не поставя: доколко конструирането по време на компилация предотвратява
инжекциите и къде спира тази защита. Приносите, формулирани в заключението, се
отнасят само към това изследване и към промените в библиотеката, произтекли от
него.

За изложението са необходими три понятия от платформата. \textbf{Свойство}
(\code{orm::property<T, "name">}) описва колона: типът на стойността и името на
колоната са аргументи на шаблона. \textbf{Междинното представяне} на заявката
(\textit{IR}) е стойност от тип като \code{select\_query<...>}, чиито параметри на
шаблона кодират избраните колони, условията и ограниченията. \textbf{Конекторът}
превежда междинното представяне към езика на конкретно хранилище и предава
стойностите на драйвера.

\subsection{Откъде идва всяка част на заявката}
\label{sec:enf-origin}

Листинг~\ref{lst:enf-query} показва заявка, написана по документирания начин.

\begin{lstlisting}[language=C++,caption={Заявка, написана по документирания начин},label={lst:enf-query}]
struct Item
{
    orm::property<int, "id">           id;    // името на колоната е аргумент на шаблон
    orm::property<std::string, "name"> name;
};

namespace orm
{
    template <>
    struct table_name_trait<Item>
    {
        static constexpr std::string_view value = "items";   // константен израз
    };
}

constexpr auto q = orm::select(orm::field<&Item::id>)
                       .where(orm::field<&Item::name> == orm::Placeholder<std::string>{});

auto rows = db.execute(q, input);   // стойността пътува отделно от заявката
\end{lstlisting}

В листинга трябва да се забележат три неща: името на колоната е част от типа на
свойството, името на таблицата е константен израз, а стойността \code{input} не
участва в изграждането на \code{q} и се подава едва при изпълнението.
Таблица~\ref{tab:enf-origin} обобщава произхода на всяка част.

\begin{table}[H]
\centering
\caption{Произход на частите на заявката}
\label{tab:enf-origin}
\begin{tabularx}{\textwidth}{L{3.6cm}L{4.6cm}Y}
\toprule
\textbf{Част на заявката} & \textbf{Произход} & \textbf{Достъпна ли е за вход по време на изпълнение} \\
\midrule
Имена на таблици и колони & \code{constexpr} метаданни на указател към член (\code{details/member\_pointer.hpp:44}, \code{entity/property.hpp:19}) & Не \\
Форма на заявката (избор, условие, съединение, подреждане) & типове в междинното представяне & Не \\
Стойности & типизирани заместители \code{Placeholder<T>} и \code{ph<T, \_N>} & Да, само като данни \\
\bottomrule
\end{tabularx}
\end{table}

Таблицата съдържа твърдение, което не следва автоматично от синтаксиса на
листинга: че няма \textit{друг} път, по който стойност да достигне заявката. Това
твърдение се оказа невярно за предходната версия на библиотеката и е обект на
казуса в глава~\ref{sec:case-study}.

\subsection{Механизми при отделните хранилища}
\label{sec:enf-backends}

Гаранцията е една и съща за всички хранилища, но механизмът, по който тя се
постига, не е. Всеки конектор използва най-силното средство, което драйверът му
предоставя. Таблица~\ref{tab:enf-mechanisms} показва резултата от анализа на
изходния код.

\begin{table}[H]
\centering
\caption{Как стойността достига драйвера при всеки конектор}
\label{tab:enf-mechanisms}
\begin{tabularx}{\textwidth}{L{2.1cm}L{4.0cm}L{2.9cm}Y}
\toprule
\textbf{Хранилище} & \textbf{Механизъм} & \textbf{Вид} & \textbf{Място в кода} \\
\midrule
PostgreSQL & \code{PQexecParams} с масиви от параметри & свързване в сървъра & \code{postgresql\_live.hpp:552, 615} \\
MySQL & подготвени заявки с \code{MYSQL\_BIND} & свързване в сървъра & \code{mysql\_live.hpp:316} \\
SQLite & \code{sqlite3\_prepare\_v2} и \code{sqlite3\_bind\_*} & свързване & \code{sqlite\_db.hpp:62, 67, 72} \\
Cassandra & \code{cass\_statement\_bind\_*} & свързване в сървъра & \code{cassandra\_live.hpp:200-221} \\
Neo4j & параметри на Bolt чрез \code{neo4j\_map\_entry} & типизирани параметри на драйвера & \code{neo4j\_live.hpp:174} \\
Redis & форми с форматиращ низ и \code{argv} на \code{hiredis} & кодиране на протоколно ниво & \code{redis\_live.hpp:109, 180, 257} \\
MongoDB & построител на документи \code{bsoncxx} & структурно, без текстов език & \code{mongodb\_live.hpp:214-238} \\
\bottomrule
\end{tabularx}
\end{table}

От таблицата се извеждат четири вида механизъм. Четири конектора използват
свързване на параметри. Neo4j използва типизирани параметри на драйвера. Redis
разчита на кодирането на протокола. MongoDB изобщо не изгражда текст. Изводът е,
че една и съща абстракция се защитава пред четири различни протокола и
аргументът за сигурност трябва да е валиден поотделно за всеки от тях.

\subsubsection{Хранилища без текстов език за заявки}

При MongoDB и Redis стойността никога не става част от текст, който хранилището
анализира синтактично. MongoDB получава документ BSON, в който стойността е
поле със собствен тип. Redis получава кадри RESP, в които всеки аргумент е
предшестван от дължината си. Липсва граматика, от която низ на нападателя би
могъл да излезе. Тази позиция е по-силна от свързването на параметри, а не
равностойна на него, тъй като не зависи от правилно разделяне на маркери и
стойности. Разликата обикновено се губи, когато инжекциите в NoSQL се разглеждат
като допълнение към SQL инжекциите.

Тази позиция е валидна само докато стойността не се интерпретира като оператор.
При MongoDB конекторът добавя стойността като типизирано поле на документа
(\code{mongodb\_live.hpp:214-238}), а не разчита произволен документ от входа.
Затова низ, който съдържа обект с оператора {\ttfamily \$ne}, остава низ.

\subsubsection{Neo4j: форматиране, което не е интерполация}

Конекторът за Neo4j изгражда текста на Cypher чрез \code{std::format}, например
шаблона \code{MATCH (n:\{\}) RETURN \{\}\{\}} (\code{neo4j\_live.hpp:261}). Това
изглежда като сглобяване на низ и заслужава внимание. Вмъкваните полета обаче са
етикетът на възела и проекцията, и двете идват от \code{constexpr} метаданни на
типа, а не от входа. Стойностите по време на изпълнение се предават отделно,
като типизирани параметри на Bolt (\code{neo4j\_live.hpp:174}). Форматирането на
константи от времето на компилация в шаблон не е интерполация на недоверени
данни. Различието е съществено, но то трябва да се демонстрира, а не да се
приема. Докъде е стигнала тази демонстрация е разгледано в
точка~\ref{sec:eval-servers}.

\subsubsection{Redis: форматиращият низ не е \texttt{printf}}

Конекторът за Redis извиква \code{redisCommand} с форматиращ низ, в който
аргументът е означен с {\ttfamily \%s} (\code{redis\_live.hpp:180}). Въпреки
сходството с \code{printf}, \code{hiredis} не замества текст. Всеки аргумент се
кодира като отделен елемент на RESP с префикс за дължина, така че съдържанието му
не може да бъде прочетено като протокол. Формата с \code{argv}
(\code{redis\_live.hpp:257}) прави това изрично.

\subsubsection{Управление на транзакции}

Командите за транзакции изпращат литералните константи \code{BEGIN},
\code{COMMIT} и \code{ROLLBACK} (\code{postgresql\_live.hpp:506, 518, 530}). Те
не съдържат потребителски данни и не нарушават гаранцията.

\subsection{Проверка на границата}
\label{sec:enf-boundary}

Анализът по хранилища потвърждава предпоставка 2 за всеки конектор поотделно, но
оставя един въпрос: стига ли всеки път към конектор само до добре формирано
междинно представяне. В хода на настоящата работа беше установено, че отговорът е
бил отрицателен (глава~\ref{sec:case-study}). Затова в библиотеката е добавена
проверка, извиквана от всяка входна точка, която достига конектор.

\begin{lstlisting}[language=C++,caption={Проверката на границата (съкратено)},label={lst:enf-shape}]
template <typename Query>
constexpr void check_query_shape()
{
    using Q = std::remove_cvref_t<Query>;
    static_assert(is_query_ir<Q>::value,
        "orm: this is not a query. ...");
    if constexpr (is_query_ir<Q>::value)
        static_assert(!query_has_literal<Q>::value,
            "orm: a literal operand reaches execution. ...");
}
\end{lstlisting}

Проверката в листинг~\ref{lst:enf-shape} има две части. Първата отхвърля всичко,
което не е един от четирите вида заявки (\code{select\_query},
\code{insert\_query}, \code{update\_query}, \code{delete\_query}). Анализът на
всички конектори показа, че това са единствените видове, които те изпълняват,
така че проверката не отхвърля нито една легитимна употреба. Втората част
обхожда условията, съединенията и присвояванията в заявката и отхвърля всеки
\textbf{литерален операнд}: стойност, записана направо в заявката, вместо през
заместител. Сравнението с \code{nullptr} и сравнението между две колони не са
литерали и остават разрешени.

Проверката се извиква от седемте входни точки, които достигат конектор:
\code{db::operator<<}, \code{db::execute}, \code{db::find\_one},
\code{db::prepare}, \code{prepared\_query::execute} и двете асинхронни
съответствия в \code{async\_db}. Пропускането на която и да е от тях би оставило
заобиколен път. Затова всяка входна точка се покрива от отделен компилационен
тест (глава~\ref{sec:evaluation}).

\subsection{Известен остатъчен риск}
\label{sec:enf-residual}

Анализът установи и един път, който е по-слаб от свързването, без да е
използван. Функцията \code{to\_sql\_literal} (\code{mysql\_live.hpp:249})
изгражда SQL чрез \textbf{екраниране} и вмъкване на стойността, вместо чрез
свързване. Коментарът ѝ я описва като предназначена за заявки без подготвени
изрази. Търсене в каталозите \code{lib} и \code{tests} намира само дефиницията и
нито едно извикване. Следователно тя е мъртъв код и не засяга текущата гаранция.
Тя е отбелязана, защото екранирането е по-слаб механизъм
(точка~\ref{sec:theory-remedies}) и бъдещо използване би отслабило гаранцията за
MySQL, без това да е видимо от интерфейса. Препоръката е функцията да бъде
премахната или изрично изолирана.

\newpage
\section{Казус: непривързаните литерални операнди}
\label{sec:case-study}

Глава~\ref{sec:enforcement} установи, че всеки конектор свързва стойностите на
заместителите правилно. Тази глава проследява случай, в който предпоставка 2 от
точка~\ref{sec:theory-compile-time} не е била изпълнена, въпреки че всеки
конектор поотделно работи коректно. Казусът е избран по две причини. Той
демонстрира метода на работата върху реален дефект, от хипотезата през
измерването до поправката. Той показва и защо анализът на всеки компонент
поотделно не е достатъчен, когато гаранцията зависи от това, какво изобщо може да
достигне компонентите.

\subsection{Откриване}
\label{sec:cs-discovery}

Отправната точка беше въпросът дали типизираният заместител е
\textit{единственият} път, по който стойност достига заявката. Операторите за
сравнение на полета са дефинирани чрез макрос в слоя на заявките.
Листинг~\ref{lst:cs-operator} показва съществената му част.

\begin{lstlisting}[language=C++,caption={Операторът за сравнение на поле (съкратено)},label={lst:cs-operator}]
template <auto P, typename T>
    requires(!std::is_same_v<T, std::nullptr_t>)
constexpr auto operator==(mem_ptr<P>, T x)
{
    static_assert(comparable_value_types<...>, "...");
    if constexpr (is_field<T>)
        return Rule<decltype(P), "==", decltype(T::get())>(P, T::get());
    else
        // стойността се записва в заявката
        return Rule<decltype(P), "==", T>(P, x);
}
\end{lstlisting}

От листинга се вижда, че параметърът \code{T} не е ограничен до заместители.
Операторът приема всяка стойност със съвместим тип и я записва в междинното
представяне като десен операнд. Следователно съществува и втори път: стойност,
записана направо в условието, например \code{field<\&Item::name> == input}, без
заместител. Такава стойност се нарича в тази работа \textbf{литерален операнд}.
Ако тя бъде вмъкната в текста на заявката, гаранцията е нарушена. Ако бъде
свързана, гаранцията е запазена. Кодът на оператора не дава отговор, тъй като
решението се взема в конектора.

\subsection{Хипотеза и измерване}
\label{sec:cs-measurement}

Конекторът за SQLite превежда всеки операнд, който не е колона и не е
заместител, в маркер \code{?} (листинг~\ref{lst:cs-render}, последният клон).
Литералът не става текст, което изключва инжекцията. Търсене на код, който след това свързва
стойността на литерала, не откри такъв в нито един от осемте конектора (седемте
работещи и \code{MockDB}). Хипотезата беше, че маркерът остава несвързан и
драйверът го тълкува като \code{NULL}.

\begin{lstlisting}[language=C++,caption={Преводът на операнд при SQLite преди поправката (съкратено)},label={lst:cs-render}]
template <typename T>
std::string render_operand_sqlite(const T& /*v*/)
{
    if constexpr (is_field<T>)
        return std::string(T::column_name());
    // ...
    else if constexpr (is_placeholder_v<T>)
    {
        // ...
    }
    else
        return "?";   // литерален операнд: маркер без стойност
}
\end{lstlisting}

Хипотезата беше проверена с две програми върху SQLite в паметта, с библиотеката
във версията преди поправката (\code{d5795a9}). Таблица~\ref{tab:cs-select}
показва резултата за заявка за избор.

\begin{table}[H]
\centering
\caption{Заявка за избор: заместител спрямо литерал}
\label{tab:cs-select}
\begin{tabularx}{\textwidth}{L{4.2cm}L{4.2cm}L{2.2cm}Y}
\toprule
\textbf{Път} & \textbf{Вход} & \textbf{Редове} & \textbf{Правилно} \\
\midrule
типизиран заместител & \code{banana} & 1 & 1 \\
типизиран заместител & \code{x' OR '1'='1} & 0 & 0 \\
литерал в заявката & \code{banana} & \textbf{0} & 1 \\
литерал в заявката & \code{x' OR '1'='1} & 0 & 0 \\
\bottomrule
\end{tabularx}
\end{table}

Резултатът потвърждава хипотезата. Пътят през заместителя връща правилния
отговор и при безобиден вход, и при полезен товар. Пътят през литерала връща
нула редове дори за \code{banana}, за който съществува съвпадащ ред. Полезният
товар не постига нищо и по двата пътя, но това не означава едно и също: при
литерала неуспехът му се дължи на това, че заявката изобщо не работи.

Таблица~\ref{tab:cs-update} показва същото при промяна и изтриване.

\begin{table}[H]
\centering
\caption{Промяна и изтриване: заместител спрямо литерал}
\label{tab:cs-update}
\begin{tabularx}{\textwidth}{L{5.4cm}L{4.4cm}Y}
\toprule
\textbf{Случай} & \textbf{Намерение} & \textbf{Резултат} \\
\midrule
заместител в \code{SET} & преименуване на ред 1 & редът е преименуван \\
литерал в \code{SET} & преименуване на ред 2 & \textbf{редът е непроменен} \\
литерал в \code{WHERE} & изтриване на ред 3 & \textbf{нищо не е изтрито} \\
\bottomrule
\end{tabularx}
\end{table}

Листинг~\ref{lst:cs-probe} показва дословния изход на втората програма.

\begin{lstlisting}[language={},deletekeywords={orm},caption={Изход на програмата за промяна и изтриване},label={lst:cs-probe}]
U1 set via placeholder 'grape', id=1   -> name now grape
U2 set via literal     'grape', id=2   -> name now banana
D1 delete via literal id==3            -> rows left 3 (3 means nothing deleted)
\end{lstlisting}

\subsection{Анализ на причината}
\label{sec:cs-cause}

Случаят с промяната заслужава отделно внимание, защото резултатът му не е
просто липса на действие. Аргументите, подадени на \code{execute()}, се свързват
позиционно, от позиция 1 нататък (листинг~\ref{lst:enf-sqlite-bind}). Заявката
\code{UPDATE items SET name = ? WHERE id = ?} има два маркера, а при изпълнението
е подаден един аргумент, предназначен за условието. Маркерът на литерала в
\code{SET} е първи. Затова аргументът е свързан с него, а маркерът в условието е
останал несвързан, тоест \code{NULL}, и не е съвпаднал с нито един ред. Това е
\textbf{объркване на параметрите}: стойност на извикващия код попада на погрешна
позиция.

От гледна точка на сигурността дефектът е \textbf{затворен отказ}: при грешка
системата остава в безопасно състояние. Литералът никога не става текст и не
може да промени структурата. Дефектът обаче не е безобиден. Той води до тихо
грешни резултати, без грешка и без предупреждение. За приложение, което
проверява правомощия чрез заявка, грешен празен резултат може да означава отказ
на легитимен достъп. Случай с видим погрешен ефект, например промяна на погрешен
ред, не беше наблюдаван в измерените заявки, но не е изключен за всяка форма на
заявка.

Остава въпросът защо дефектът не е бил открит от съществуващите тестове. Търсене
в тестовете намери литерални операнди само в единичните тестове на междинното
представяне, които проверяват формата на междинното представяне чрез
\code{static\_assert} и никога не изпълняват заявката. Нито един тест не е
изпълнявал заявка с литерален операнд срещу база данни. Дефектът е оцелял в
сляпо петно на тестовото покритие, а не поради проектно решение.

\subsection{Избор на поправка}
\label{sec:cs-fix}

Възможни са две поправки.

\textbf{Свързване на литералите.} Всеки конектор събира стойностите на
литералите и ги свързва заедно с аргументите на \code{execute()}, в общ ред на
позициите. Така литералите стават законен втори път за стойности.

\textbf{Отхвърляне на литералите.} Заявка с литерален операнд не се компилира, ако
бъде подадена за изпълнение. Единственият път за стойности остава заместителят.

Изборът падна върху отхвърлянето по три причини. Първо, свързването изисква
промяна в превода и свързването на всички седем конектора. Грешка в реда на
позициите в който и да е от тях би превърнала затворения отказ в объркване на
параметрите с видим ефект, а реализация, която вмъква литерала като текст, би
създала истинска инжекция. Отхвърлянето е една проверка на едно място
(точка~\ref{sec:enf-boundary}). Второ, то съответства на тезата на работата:
стойностите влизат само през заместители. Трето, поправката не чупи работещ код,
защото кодът, който тя отхвърля, и без това е давал грешни резултати.

Цената е, че естествени записи, например условие за положително число, вече
изискват заместител. Отхвърлянето засяга и литерали, известни по време на
компилация, защото дефектът никога не е зависел от това кога е известна
стойността: нито един конектор не свързва литерал.

\subsection{Проверка на поправката}
\label{sec:cs-verification}

Поправката е проверена в два етапа. Преди нея наборът от компилационни тестове
беше изпълнен с библиотеката във версия \code{d5795a9}. Седемте теста за
литерални операнди, по един за всяка входна точка и вид заявка, се провалиха с
един и същ резултат: програмата се компилира, въпреки че не би трябвало.
От 21 теста в това изпълнение 12 преминаха, а 9 се провалиха: седемте за
литерални операнди и двата, разгледани в точка~\ref{sec:eval-compile}.

След поправката и седемте програми се отхвърлят при компилация по проверяваната
причина. Всички тестове преминават: 314-те съществуващи в библиотеката и 23-те
нови от набора за оценка. Последното е съществено: то показва, че поправката
затваря дефекта, без да засегне работещ код. Подробните резултати са в
глава~\ref{sec:evaluation}.

\subsection{Изводи от казуса}
\label{sec:cs-lessons}

Казусът води до два извода, които надхвърлят конкретния дефект. Първият е, че
гаранцията за сигурност на една абстракция не е сума от гаранциите на
компонентите ѝ. Всеки конектор свързва заместителите правилно, но никой от тях не
е отговарял за въпроса какво изобщо може да ги достигне. Вторият е, че тест,
който проверява само формата на заявката, не проверява поведението ѝ. Дефектът
оцеля, защото литералите бяха тествани като структура, но не и като изпълнени
заявки. Методът на настоящата работа, който изисква всяко твърдение да бъде
измерено, а не изведено, е причината дефектът да бъде открит.

\newpage
\section{Оценка}
\label{sec:evaluation}

Главите~\ref{sec:enforcement} и~\ref{sec:case-study} описаха как гаранцията се
постига и как беше затворен един пропуск в нея. Настоящата глава проверява
гаранцията по метода от точка~\ref{sec:tm-method} и съпоставя резултатите с
критериите от таблица~\ref{tab:tm-criteria}. Наблюденията се представят първо, а
тълкуването им следва отделно, за да може читателят да прецени едното независимо
от другото.

\subsection{Среда и възпроизводимост}
\label{sec:eval-env}

Всички измервания са извършени с изходния код във версия \code{a17f4e7}, при
следната среда:

\begin{itemize}
    \item операционна система Ubuntu 24.04.4 LTS под WSL2, ядро
          \code{6.6.114.1-microsoft-standard-WSL2};
    \item компилатори GCC 14.2.0 и Clang 18.1.3, стандарт C++23;
    \item SQLite 3.45.1, база данни в паметта, нова за всяко изпълнение;
    \item CMake 3.28.3.
\end{itemize}

Наборът се изгражда и изпълнява с опцията \code{ORM\_ENABLE\_SECURITY}, както е
описано в приложение~А. Всеки резултат се записва в каталога
\code{security/evidence/} в изходния код, а не в каталога за изграждане. Така
всяко изпълнение е преглеждаема разлика в системата за контрол на версиите и
всяко число в тази глава може да бъде проследено до файл. От версия
\code{a17f4e7} нататък наборът се изпълнява и при всяка промяна в системата за
непрекъсната интеграция, като изпълнението за самата версия \code{a17f4e7} е
успешно.

\subsection{Корпус от полезни товари}
\label{sec:eval-corpus}

Корпусът съдържа 66 полезни товара от 11 публични източника. Всеки ред съдържа
препратка към източника си и нито един полезен товар не е съставен за целите на
работата. Таблица~\ref{tab:eval-corpus} показва разпределението.

\begin{table}[H]
\centering
\caption{Състав на корпуса}
\label{tab:eval-corpus}
\begin{tabularx}{\textwidth}{L{3.0cm}L{2.0cm}Y}
\toprule
\textbf{Клас} & \textbf{Брой} & \textbf{Диалекти} \\
\midrule
SQL & 36 & общи 16, MySQL 8, PostgreSQL 6, SQLite 6 \\
NoSQL (MongoDB) & 15 & инжекция на оператори \\
Cypher (Neo4j) & 6 & \\
CQL (Cassandra) & 4 & \\
Redis & 5 & инжекция на протоколно ниво \\
\midrule
Общо & 66 & \\
\bottomrule
\end{tabularx}
\end{table}

Разпределението е неравномерно и това е следствие от правилото полезен товар да
се включва само ако е открит дословно в източник. Три полезни товара са
публикувани от източника с незапълнено поле, например \code{[RANDNUM]}. Те се
изпълняват точно както са публикувани, защото запълването на полето би означавало
съставяне на нов текст. Корпусът изрично изброява и пропуските си, сред които
липсата на проверен полезен товар с \code{UNION} и с натрупани заявки за SQLite,
както и с \code{UNION} за Cypher.

\subsection{Отрицателни компилационни тестове}
\label{sec:eval-compile}

Наборът съдържа 18 компилационни теста: 15 програми, които трябва да бъдат
отхвърлени, и 3 контролни случая, които трябва да се компилират.
Таблица~\ref{tab:eval-compile} показва резултата при двата компилатора.

\begin{table}[H]
\centering
\caption{Резултати от компилационните тестове}
\label{tab:eval-compile}
\begin{tabularx}{\textwidth}{L{6.3cm}L{1.6cm}L{1.9cm}YY}
\toprule
\textbf{Случай} & \textbf{Ниво} & \textbf{Очакване} & \textbf{GCC} & \textbf{Clang} \\
\midrule
00 документирана заявка        & контрола & компилира & да & да \\
14 сравнение с \code{nullptr}  & контрола & компилира & да & да \\
15 сравнение на две колони      & контрола & компилира & да & да \\
\midrule
06 низ към целочислен заместител & 1 & отхвърля & да & да \\
01 име на колона по време на изпълнение & 2 & отхвърля & да & да \\
16 име на колона в \code{std::string} & 2 & отхвърля & да & да \\
02 име на таблица по време на изпълнение & 2 & отхвърля & да & да \\
04 колона, избрана по време на изпълнение & 2 & отхвърля & да & да \\
17 две колони в една променлива & 2 & отхвърля & да & да \\
03 текст на заявка към \code{execute} & 3 & отхвърля & да & да \\
05 условие като текст & 3 & отхвърля & да & да \\
\midrule
07 литерал, \code{execute} & литерал & отхвърля & да & да \\
08 литерал, \code{operator<<} & литерал & отхвърля & да & да \\
09 литерал, \code{find\_one} & литерал & отхвърля & да & да \\
10 литерал, \code{prepare} & литерал & отхвърля & да & да \\
11 литерал в \code{SET} & литерал & отхвърля & да & да \\
12 литерал в \code{DELETE} & литерал & отхвърля & да & да \\
13 литерал, константа & литерал & отхвърля & да & да \\
\bottomrule
\end{tabularx}
\end{table}

„Да“ в таблицата означава, че програмата е дала очаквания резултат и, при
отхвърляне, че диагностиката съответства на образеца за проверяваната причина.
Диагностиките са запазени поотделно за всеки компилатор в
\code{security/evidence/compile/}.

Два от тестовете бяха уточнени след първото изпълнение и това има значение за
тълкуването. Първоначалните варианти на случаи 01 и 04 бяха отхвърлени, но по
причина, различна от проверяваната, и затова се отчетоха като неуспешни. Случай
01 опитваше да използва \code{std::string} като име на колона и беше отхвърлен
поради несъответствие на типа, преди въпросът за константността изобщо да
възникне. Той беше разделен на случай 01, който подава стойност с правилния тип,
но известна едва по време на изпълнение, и случай 16, който запазва варианта с
\code{std::string}. Случай 04 опитваше да избере между две колони и беше
отхвърлен със съобщението, че двата указателя към член имат различни типове.
Това разкри по-силно свойство от проверяваното: понеже името на колоната е част
от типа на свойството, две колони с един и същ тип на стойността (тук и двете
\code{int}) са различни типове, и една променлива не може да обхваща и двете.
То е отделено в случай 17, а случай 04 проверява само константността.
Записът от първото изпълнение е запазен в
\code{security/evidence/compile\_before\_fix.txt}.

\subsection{Диференциален тест със SQLite}
\label{sec:eval-differential}

От корпуса към SQLite са приложими 22 полезни товара: общите 16 и специфичните
за SQLite 6. Таблица~\ref{tab:eval-diff-summary} обобщава резултата.

\begin{table}[H]
\centering
\caption{Диференциален тест със SQLite, 22 полезни товара}
\label{tab:eval-diff-summary}
\begin{tabularx}{\textwidth}{L{5.6cm}YY}
\toprule
\textbf{Резултат} & \textbf{Наивна заявка} & \textbf{Библиотеката} \\
\midrule
Компрометирана: върнати редове & 4 & 0 \\
Компрометирана: промяна на състоянието & 0 & 0 \\
Компрометирана: принудително изпълнение & 0 & 0 \\
Синтактично нарушена & 10 & 0 \\
Без наблюдаем ефект & 8 & 22 \\
\bottomrule
\end{tabularx}
\end{table}

Таблица~\ref{tab:eval-diff-technique} показва резултата на наивната заявка по
техники.

\begin{table}[H]
\centering
\caption{Резултат на наивната заявка по техники}
\label{tab:eval-diff-technique}
\begin{tabularx}{\textwidth}{L{4.2cm}L{1.6cm}Y}
\toprule
\textbf{Техника} & \textbf{Брой} & \textbf{Резултат} \\
\midrule
тавтология & 2 & и двата компрометират (3 и 1 върнати реда) \\
\code{UNION} & 2 & и двата компрометират (по 1 върнат ред) \\
сляпа атака по време & 6 & и шестте нарушават синтаксиса \\
натрупана заявка & 1 & нарушава синтаксиса \\
сляпа атака по истинност & 3 & 2 без ефект, 1 нарушава синтаксиса \\
коментар & 3 & 2 без ефект, 1 нарушава синтаксиса \\
изброяване & 2 & 1 без ефект, 1 нарушава синтаксиса \\
предизвикване на грешка & 1 & без ефект \\
заобикаляне на екраниране & 2 & без ефект \\
\bottomrule
\end{tabularx}
\end{table}

Наблюденията са следните. Наивната заявка е компрометирана от 4 от 22 полезни
товара. И четирите са върнали редове, които заявката не би трябвало да върне.
Единият от двата полезни товара с \code{UNION}, \code{SQL-003}, е върнал реда с
потребителското име и паролата от таблицата \code{users}, тоест е извлякъл
съхранена идентификационна информация. Другият, \code{SQL-004}, е върнал ред с
избрани от нападателя стойности, техника за заобикаляне на удостоверяване. През
библиотеката и 22-та полезни товара не са имали наблюдаем ефект. Пълният резултат
по полезни товари е в \code{security/evidence/differential\_sqlite.csv}.

\subsection{Тълкуване}
\label{sec:eval-interpretation}

Резултатът на библиотеката, 22 от 22 без ефект, придобива смисъл едва на фона на
резултата на наивната заявка. Четирите компрометирания показват, че корпусът и
постановката са способни да разкрият уязвимост, тоест че нулата при библиотеката
не се дължи на безсилен корпус.

Десетте синтактично нарушени заявки заслужават отделен коментар. Те не са
неуспех на атаката в смисъла на защитата. Полезният товар е достигнал граматиката
на SQLite и е нарушил заявката, но не е постигнал ефект. Причината е проверена за
всеки случай. Пет от шестте слепи атаки по време използват конструкции на други
системи, които SQLite не поддържа: \code{SLEEP} на MySQL, \code{WAITFOR DELAY} на
SQL Server и \code{dbms\_pipe} на Oracle. Шестата, \code{SQL-034}, е специфична за
SQLite, но е публикувана с незапълнени полета и затова не е синтактично валидна.
Натрупаната заявка \code{SQL-017} извиква процедурата \code{xp\_cmdshell} на SQL
Server. Срещу библиотеката същите полезни товари не достигат граматиката изобщо:
те пристигат като стойност на свързан параметър и се сравняват като низ.

Осемте полезни товара без ефект при наивната заявка са предимно слепи атаки и
опити за заобикаляне на екраниране. Слепата атака по дефиниция не дава видим
резултат при едно изпълнение: тя разчита на разлика между две изпълнения. Затова
оракулът, който вижда само ефекти, не я отчита. Както е посочено в
точка~\ref{sec:tm-not-proven}, броят на компрометиранията на наивната заявка е
долна граница.

\subsection{Сървърните хранилища}
\label{sec:eval-servers}

Диференциалният тест е изпълнен със SQLite. За останалите шест хранилища
приложимите полезни товари са: PostgreSQL 22, MySQL 24, MongoDB 15, Neo4j 6,
Cassandra 4 и Redis 5. Изпълнението им изисква работещи сървъри, които се
стартират в контейнери. В средата, в която са извършени измерванията за
настоящата версия на работата, услугата за контейнери не беше налична и тези
изпълнения не са извършени. За тези хранилища гаранцията е потвърдена чрез
анализа на изходния код в точка~\ref{sec:enf-backends}, но не и чрез измерване.

Конекторът за Neo4j има допълнително ограничение. Образът за контейнери, използван
в системата за непрекъсната интеграция, не съдържа клиентската библиотека
\code{libneo4j-client} и не стартира сървър Neo4j. Затова за Neo4j и
съществуващите интеграционни тестове на библиотеката не се изпълняват. Измерването
му изисква разширяване на образа и е посочено като бъдеща работа.

\subsection{Съпоставка с критериите за успех}
\label{sec:eval-criteria}

Таблица~\ref{tab:eval-criteria} съпоставя резултатите с критериите от
таблица~\ref{tab:tm-criteria}.

\begin{table}[H]
\centering
\caption{Съпоставка с критериите за успех}
\label{tab:eval-criteria}
\begin{tabularx}{\textwidth}{L{1.0cm}L{2.8cm}Y}
\toprule
\textbf{№} & \textbf{Оценка} & \textbf{Основание} \\
\midrule
К1 & изпълнен за SQLite & 22 от 22 полезни товара без ефект през библиотеката; за сървърните хранилища само чрез анализ \\
К2 & изпълнен & 15 от 15 отхвърляния по проверяваната причина и 3 от 3 контролни случая, при GCC и при Clang \\
К3 & изпълнен за SQLite & наивната заявка е компрометирана от 4 полезни товара \\
К4 & изпълнен & назовани граници (точка~\ref{sec:tm-scope}), демонстриран и затворен пропуск (глава~\ref{sec:case-study}), отбелязан остатъчен риск (точка~\ref{sec:enf-residual}) \\
\bottomrule
\end{tabularx}
\end{table}

Критериите К2 и К4 са изпълнени изцяло. Критериите К1 и К3 са изпълнени за
SQLite, а за сървърните хранилища К1 се подкрепя само от анализа на изходния код.
Разликата между двата вида подкрепа е съществена и не се замъглява: анализът
показва механизма, но само измерването може да покаже, че механизмът работи
така, както е описан.

\subsection{Заплахи за валидността}
\label{sec:eval-validity}

\textbf{Вътрешна валидност.} Оракулът вижда само ефекти и подценява наивния път.
Това изместване работи срещу тезата, а не в нейна полза. Компилационните тестове
използват конектора \code{MockDB}, който превежда заявките без база данни.
Проверката на границата обаче се извършва в общия слой, преди конектора, така че
изборът на конектор не влияе на резултата от тези тестове.

\textbf{Конструктивна валидност.} Диференциалният тест използва една форма на
заявка с две колони. Тя е избрана, за да дава шанс на най-разпространените
полезни товари с \code{UNION}. Друга форма би могла да направи наивния път
уязвим за други полезни товари. Това не засяга резултата на библиотеката,
защото при нея стойността не достига граматиката при никоя форма.

\textbf{Външна валидност.} Измерването е извършено с един диалект и един корпус
от 66 полезни товара. Корпусът не е изчерпателен и съдържа пропуски, които са
изброени в самия него. Твърдението, че гаранцията се отнася за всички входове, не
се основава на корпуса, а на аргумента от точка~\ref{sec:theory-compile-time},
който компилационните тестове проверяват. Корпусът показва, че аргументът издържа
на реални атаки, но не замества аргумента.

\textbf{Надеждност.} Резултатите са детерминирани. Измерванията на времето са
изключени от записаните резултати, за да не предизвикват разлики между
изпълненията. Всеки резултат е записан във файл в изходния код и може да бъде
получен повторно с командите от приложение~А.

\newpage
\section*{ЗАКЛЮЧЕНИЕ}
\label{sec:conclusion}
\addcontentsline{toc}{section}{ЗАКЛЮЧЕНИЕ}

Работата постави въпроса дали една абстракция, конструираща заявките по време на
компилация, може да елиминира цяло семейство от класове инжекции при хетерогенни
хранилища за данни, и къде спира тази гаранция.

\subsection*{Отговор на изследователския въпрос}

Отговорът е положителен, с ясно назовани граници. Когато имената на таблици и
колони и формата на заявката са аргументи на шаблони, а стойностите достигат
заявката само чрез свързани типизирани заместители, недоверен вход не може да
промени структурата на заявката. Това следва от определението на инжекцията като
промяна на синтактичната структура~\cite{su2006essence} и от изискването
аргументите на шаблон да бъдат константни изрази. Аргументът е проверен, а не
само изведен: 15 от 15 опита идентификатор или структура да зависят от стойност
по време на изпълнение са отхвърлени от компилатора по проверяваната причина, при
два компилатора, а при SQLite нито един от 22 приложими полезни товара не е имал
ефект през библиотеката, докато наивната заявка е компрометирана от 4.

Гаранцията обаче не беше пълна в началото на изследването. Стойност, записана
направо в заявката, се приемаше и оставаше несвързана. Това не позволяваше
инжекция, но водеше до тихо грешни резултати. Пропускът е затворен. Именно този
резултат показва най-ясно защо втората част на изследователския въпрос е
равностойна на първата: гаранцията беше изказана като пълна, преди да бъде
измерена.

\subsection*{Приноси}

\begin{enumerate}
    \item \textbf{Научно-приложен принос.} Аргумент, който свързва формалното
          определение на инжекцията с изискването за константни изрази в C++ и
          го разширява от SQL към пет класа инжекции при седем хранилища.
          Установено е, че гаранцията се постига с четири различни механизма и
          че при две от хранилищата липсва текстова граматика, от която полезен
          товар би могъл да излезе.
    \item \textbf{Приложен принос.} Възпроизводим набор за оценка, включен в
          системата за непрекъсната интеграция на библиотеката: 18 компилационни
          теста с контролни случаи и проверка на диагностиката, корпус от 66
          полезни товара с посочен източник за всеки и диференциален тест със
          съзнателно уязвима база за сравнение.
    \item \textbf{Приложен принос.} Открит, измерен и затворен пропуск в
          гаранцията на библиотеката, както и проверка на границата на
          интерфейса, която отхвърля текст на заявка със съобщение, назоваващо
          причината.
    \item \textbf{Отрицателни резултати.} Назовани граници на гаранцията и
          остатъчен риск в неизползван път с по-слаб механизъм, който е
          мъртъв код, но би отслабил гаранцията, ако бъде използван.
\end{enumerate}

\subsection*{Ограничения}

Основното ограничение е, че диференциалният тест е изпълнен само със SQLite. За
останалите шест хранилища гаранцията е потвърдена чрез анализ на изходния код, но
не и чрез измерване. Конекторът за Neo4j не се изпълнява дори от съществуващите
интеграционни тестове на библиотеката. Оракулът вижда само ефекти и не отчита
слепите атаки. Тестът използва една форма на заявка, а корпусът не е изчерпателен.
Най-сетне, самата гаранция има цена: схеми, определяни по време на изпълнение, не
се поддържат, защото именно невъзможността идентификатор да зависи от входа е
същността на защитата.

\subsection*{Бъдеща работа}

Насоките произтичат пряко от ограниченията. Първата е изпълнението на
диференциалния тест при сървърните хранилища. Системата за непрекъсната
интеграция на библиотеката вече стартира PostgreSQL, MySQL, MongoDB, Redis и
Cassandra в контейнери, така че разширяването на теста към тях не изисква нова
инфраструктура. Втората е добавянето на Neo4j към образа за контейнери. Третата
е оракул по двойки, който сравнява две изпълнения и така отчита слепите атаки.
Четвъртата е премахването на мъртвия път с екраниране. Петата е алтернативният
дизайн, при който литералите се свързват вместо да се отхвърлят, оценен по
същия метод. Шестата е самостоятелен анализ на операциите по схемата и на
инжекцията от втори ред, които останаха извън обхвата.

\newpage

\bibliographystyle{ieeetr}
\bibliography{references}

\newpage
\section*{ПРИЛОЖЕНИЯ}
\label{sec:appendices}
\addcontentsline{toc}{section}{ПРИЛОЖЕНИЯ}

\subsection*{Приложение А. Възпроизвеждане на резултатите}
\label{app:reproduce}

Изходният код е публично достъпен в \url{https://github.com/cpp-for-everything/ORM}.
Резултатите със SQLite и компилационните тестове в глава~\ref{sec:evaluation} са
получени с версия \code{a17f4e7}. Резултатите със сървърните хранилища са от
изпълнение \code{41870491654} на работния процес \code{Docker Integration Tests}
върху комит \code{61255de}.
Наборът за оценка се изгражда и изпълнява със следните команди, от корена на
хранилището:

\begin{lstlisting}[language=bash,caption={Изграждане и изпълнение на набора за оценка},label={lst:app-reproduce}]
cmake -S . -B build-security -DORM_ENABLE_SECURITY=ON -DORM_ENABLE_SQLITE=ON
cmake --build build-security
ctest --test-dir build-security -L security --output-on-failure
\end{lstlisting}

Компилационните тестове изискват GCC или Clang. За изпълнение с Clang
променливите на средата \code{CC} и \code{CXX} се задават съответно на
\code{clang-18} и \code{clang++-18} преди конфигурирането. Резултатите се записват
в хранилището, заедно с изходния код на набора.

Програмите за сървърните хранилища се изграждат, когато съответният конектор е
включен, например с \code{-DORM\_ENABLE\_POSTGRESQL=ON}, и изискват работещ сървър,
достъпен през променливите \code{ORM\_*} на интеграционните тестове. Най-лесният
начин е образът \code{test-all} с услугите на Docker Compose, както го изпълнява
работният процес \code{Docker Integration Tests}. Той качва записаните резултати
като артефакт \code{security-evidence}.

\subsection*{Приложение Б. Състав на набора за оценка}

\begin{table}[H]
\centering
\caption{Части на набора за оценка}
\label{tab:app-files}
\begin{tabularx}{\textwidth}{L{5.4cm}Y}
\toprule
\textbf{Част} & \textbf{Съдържание} \\
\midrule
компилационни тестове & 18 програми и общата им постановка \\
проверка на компилационните тестове & проверка на изхода и на диагностиката на всеки компилационен тест \\
корпус & 66 полезни товара с колони за идентификатор, клас, диалект, техника, текст, очакван ефект и източник \\
източници на корпуса & 11 публични източника \\
тест със SQLite & диференциалният тест със SQLite \\
тестове със сървъри & пет програми: диференциалните тестове със сървърите на PostgreSQL, MySQL, MongoDB, Redis и Cassandra \\
общ драйвер & общият драйвер на сървърните тестове: контрол, оракул, резюме и правило за неуспех \\
именувани тестове & четири именувани теста на пътя през заместител \\
програми за възпроизвеждане & двете програми, възпроизвеждащи пропуска преди поправката \\
записани резултати & резюмета, записи по полезни товари и диагностики на компилаторите \\
\bottomrule
\end{tabularx}
\end{table}

Листинг~\ref{lst:app-corpus} показва вида на корпуса.

\begin{lstlisting}[language={},deletekeywords={orm},caption={Заглавен ред и първи запис на корпуса (записът е пренесен)},label={lst:app-corpus}]
"id","class","dialect","technique","payload","expected_effect","source_id"
"SQL-001","sql","generic","tautology","' OR '1'='1'--",
    "bypasses WHERE clause, returns all rows","patt_sqli"
\end{lstlisting}

\subsection*{Приложение В. Терминологичен речник}

Таблица~\ref{tab:app-glossary} съпоставя термините, използвани в текста, с
английските им съответствия и с идентификаторите в изходния код.

\begin{table}[H]
\centering
\caption{Терминологичен речник}
\label{tab:app-glossary}
\begin{tabularx}{\textwidth}{L{4.4cm}L{4.4cm}Y}
\toprule
\textbf{Термин} & \textbf{На английски} & \textbf{В кода} \\
\midrule
инжекция & injection & \\
полезен товар & payload & колона \code{payload} в корпуса \\
екраниране & escaping & \code{to\_sql\_literal} \\
свързване на параметри & parameter binding & \code{sqlite3\_bind\_*}, \code{PQexecParams} \\
типизиран заместител & typed placeholder & \code{Placeholder<T>}, \code{ph<T, \_N>} \\
литерален операнд & literal operand & \code{is\_literal\_operand\_v} \\
междинно представяне & intermediate representation & \code{select\_query} и сродните \\
свойство & property & \code{property<T, "name">} \\
конектор & connector & \code{connector\_trait<DB>} \\
проверка на границата & boundary check & \code{check\_query\_shape} \\
контролен случай & control case & случаи 00, 14 и 15 \\
оракул & oracle & функцията \code{classify} \\
затворен отказ & fail closed & \\
объркване на параметрите & parameter confusion & \\
\bottomrule
\end{tabularx}
\end{table}


\end{document}
